% Options for packages loaded elsewhere
\PassOptionsToPackage{unicode}{hyperref}
\PassOptionsToPackage{hyphens}{url}
\documentclass[
]{article}
\usepackage{xcolor}
\usepackage{amsmath,amssymb}
\setcounter{secnumdepth}{-\maxdimen} % remove section numbering
\usepackage{iftex}
\ifPDFTeX
  \usepackage[T1]{fontenc}
  \usepackage[utf8]{inputenc}
  \usepackage{textcomp} % provide euro and other symbols
\else % if luatex or xetex
  \usepackage{unicode-math} % this also loads fontspec
  \defaultfontfeatures{Scale=MatchLowercase}
  \defaultfontfeatures[\rmfamily]{Ligatures=TeX,Scale=1}
\fi
\usepackage{lmodern}
\ifPDFTeX\else
  % xetex/luatex font selection
\fi
% Use upquote if available, for straight quotes in verbatim environments
\IfFileExists{upquote.sty}{\usepackage{upquote}}{}
\IfFileExists{microtype.sty}{% use microtype if available
  \usepackage[]{microtype}
  \UseMicrotypeSet[protrusion]{basicmath} % disable protrusion for tt fonts
}{}
\makeatletter
\@ifundefined{KOMAClassName}{% if non-KOMA class
  \IfFileExists{parskip.sty}{%
    \usepackage{parskip}
  }{% else
    \setlength{\parindent}{0pt}
    \setlength{\parskip}{6pt plus 2pt minus 1pt}}
}{% if KOMA class
  \KOMAoptions{parskip=half}}
\makeatother
\usepackage{longtable,booktabs,array}
\newcounter{none} % for unnumbered tables
\usepackage{calc} % for calculating minipage widths
% Correct order of tables after \paragraph or \subparagraph
\usepackage{etoolbox}
\makeatletter
\patchcmd\longtable{\par}{\if@noskipsec\mbox{}\fi\par}{}{}
\makeatother
% Allow footnotes in longtable head/foot
\IfFileExists{footnotehyper.sty}{\usepackage{footnotehyper}}{\usepackage{footnote}}
\makesavenoteenv{longtable}
\setlength{\emergencystretch}{3em} % prevent overfull lines
\providecommand{\tightlist}{%
  \setlength{\itemsep}{0pt}\setlength{\parskip}{0pt}}
\usepackage{bookmark}
\IfFileExists{xurl.sty}{\usepackage{xurl}}{} % add URL line breaks if available
\urlstyle{same}
\hypersetup{
  hidelinks,
  pdfcreator={LaTeX via pandoc}}

\author{}
\date{}

\AtBeginDocument{\renewcommand{\setminus}{\mathbin{\backslash}}}
\begin{document}

\section{The Structural Theory of Intelligence Under Bounded
Context}\label{the-structural-theory-of-intelligence-under-bounded-context}

Patrick McCarthy

\begin{center}\rule{0.5\linewidth}{0.5pt}\end{center}

\subsection{Abstract}\label{abstract}

A companion paper {[}McCarthy2026a{]} establishes that any system
accumulating propositions under bounded active context \(C_n\) must
maintain a directed graph with typed edges --- graph structure is not a
design choice but a necessary consequence of information preservation.
This paper derives the complete structural theory of intelligence that
follows from that foundation.

Starting from bounded \(C_n\), a survival threshold \(U_{lethal}^d\),
and a world with positive entropy, we prove a connected chain of
necessary consequences: knowledge and action are inseparable (\(K\) is
constitutively indexed by the informed action space \(F\), with an
asymmetric retention criterion derived from survival pressure); the
escalation architecture dominates top-down allocation (bottom-up
failure signaling is \(O(1)\) in steady state; precision allocation
degrades as \(O(|F|/C_n)\)); uncertainty reduction follows
evaluation-driven state revision (gradient descent on \(U\) with
convergence guaranteed by a scalar Lyapunov argument); the encoding
hierarchy follows from cost structure (proven actions promote to lower
levels with rigor threshold \(\theta_i = 1 - \varepsilon/C_i\));
\(F\) and \(M\) are type-irreducible; and agency \(\mathcal{A} > 0\)
and learning efficiency \(\eta_M > 0\) are the unique structural
invariants of any persisting entity.

Intelligence is therefore the \textbf{driven capacity to improve}:
\(I(E) = \mathcal{A} \cdot \eta_M\) --- not what has been learned, but
how efficiently evidence converts into better \((F, K)\). This is not
stipulated but derived: \(\mathcal{A}\) and \(\eta_M\) are the only
properties that every persisting entity must have, and their product is
the unique normalized combination of a probability and a conditional
rate. Extensions to collective intelligence, sentience, and
superintelligence are developed in a separate companion paper.

\begin{center}\rule{0.5\linewidth}{0.5pt}\end{center}

\subsection{1. Introduction}\label{introduction}

\textbf{From graph necessity to the structural theory.} A companion
paper {[}McCarthy2026a{]} proves that any system accumulating
propositions under bounded active context must maintain a directed graph
with typed edges over its propositions. The proof is
substrate-independent: it holds for biological neural networks,
artificial systems, and institutional knowledge alike. The present paper
takes that result as a foundation and asks: \emph{given that graph
structure is necessary, what is the complete set of structural
consequences for intelligence?}

The answer is a connected chain. Starting from bounded context \(C_n\),
a survival threshold \(U_{lethal}^d\), and a world with positive
entropy, we derive: the inseparability of knowledge and action (Theorem
1); the dominance of escalation over top-down allocation (Theorems 2,
2a); evaluation-driven state revision with guaranteed convergence
(Theorem 3); the encoding hierarchy with rigor thresholds (Theorem 4);
the type irreducibility of action and learning (Theorem 5); and the
structural necessity of agency and learning efficiency (Theorem 6). Each
result is a consequence of the same generative constraint. The
individual results formalize known intuitions; the contribution is that
they form a single deductive chain from one constraint.

\textbf{The engineering origin.} This framework emerged from a concrete
engineering problem: generating repeatable, transposable tasks across a
diverse codebase ecosystem spanning thousands of repositories. The
formal theory describes what a self-optimizing agentic system converged
to under real performance constraints. The biological and civilizational
correspondences are predictions the framework makes --- not the premises
it rests on.

Provenance was the root constraint. An agentic system that synthesizes
its behavior from a context window is inherently unpredictable.
Transposability requires a quantifiable basis: evidence that an action
has worked, attributed to a source, grounded in the conditions under
which it succeeded. Without provenance --- without knowing \emph{why}
something worked --- you cannot measure confidence, and without a
confidence measure you cannot optimize. The system converged on a
knowledge graph separating \emph{why you know something} from
\emph{how to apply it}: the graph structure whose necessity is formally
established in {[}McCarthy2026a{]}.

The engineering consequence was precisely measurable. Before separating
\(K\) from \(F\): \(O(N \times M)\) retrieval cost with no sharing of
common structure. After separation: \(O(1)\) local graph traversal per
entity, independent and unbounded. The \(O(N \times M) \to O(1)\)
transformation is the empirical measurement of \(K\) precipitating from
\(F\) as the normalization of shared propositions (Corollary 2b). Adding
agents introduced the escalation architecture: subproblems that
previously required explicit orchestration became self-contained agents
that escalated only on failure. Proven subnodes converged toward
determinism and were scripted --- promoted to a lower encoding level.
The active inference load stayed bounded as the system's coverage of
\(W\) expanded.

\textbf{Biology as feasibility test.} The test of whether this is a
theory of software optimization or a theory of intelligence is whether
the formal structure is consistent with biology --- the only system
where intelligence is unambiguously instantiated over evolutionary time.
We claim that if the theory is correct, the biological architecture
should be feasible under it: that genetic encoding of \(M\) rather than
\((K, F)\) is the predicted transmission mode (Section 11, Open Question
8); that the encoding hierarchy from reflex to active reasoning follows
from the cost structure of Theorem 4; that cultural transmission
emerges as a structural necessity when the heritable channel cannot
carry individual \((K, F)\) across generations. None of these were
assumed. They follow from the same optimization constraints that
produced the software architecture.

\textbf{Relationship to prior accounts.} The framework defines a
structural class: the set of all processes satisfying bounded active
context, a survival threshold, and a growing reachable space. Biological
evolution is one member; gradient descent in artificial systems is
another; collective human knowledge transmission is another. The reason
they all converge on graph-structured knowledge, escalation
architectures, and provenance requirements is that these are necessary
properties of any member of the class (Theorem 6 and
{[}McCarthy2026a{]}).

\begin{center}\rule{0.5\linewidth}{0.5pt}\end{center}

\subsection{2. The Central Claims}\label{the-central-claims}

We claim: \textbf{intelligence is the driven capacity to improve} ---
agency \(\mathcal{A}\) exercised through a higher-order function space
\(M\) that continuously refines both the knowledge graph \(K\) and the
informed action space \(F\). Formally:
\(I(E) = \mathcal{A} \cdot \eta_M\).

Three components are jointly necessary and none is sufficient alone:

\begin{itemize}
\tightlist
\item
  \textbf{Drive (\(\mathcal{A}\)):} the intrinsic orientation toward
  uncertainty reduction --- the desire to know --- without which \(M\)
  has no reason to engage the gradient
\item
  \textbf{Mechanism (\(\eta_M\)):} the efficiency of \(M\): how
  effectively it converts gradient signals into improved \((F, K)\).
  Without \(\eta_M > 0\), the entity cannot learn from evidence; it
  executes a fixed \((F, K)\) no matter how motivated
\item
  \textbf{Outcome:} the bounding of uncertainty through
  knowledge-grounded action --- what \(I(E)\) produces over time, not
  what it is. The size of \(K\) or \(F\) measures past intelligence;
  \(I(E)\) measures the rate at which that past is being extended
\end{itemize}

\(\mathcal{A} > 0\) and \(\eta_M > 0\) are the only universal invariants
of intelligence (Theorem 6). Every other aspect of \(E\) --- \(F\),
\(K\), substrate, sensing function --- can differ arbitrarily.
\(\mathcal{A} = 0\) or \(\eta_M = 0\) yields the automaton: executing
what prior gradient descent encoded, but not itself descending.

\textbf{The contribution.} Prior theories identify individual aspects of
intelligent systems: FEP derives uncertainty minimization from
self-organization {[}Friston2010{]}; Hutter defines optimal behavior
over computable environments {[}Hutter2005{]}; Chollet measures
skill-acquisition efficiency {[}Chollet2019{]}; Simon establishes
bounded rationality {[}Simon1955{]}; Dupoux, LeCun, and Malik
{[}DupouxLeCunMalik2026{]} propose a dual-system architecture motivated
by cognitive science. Each takes its structural assumptions as given.

A companion paper {[}McCarthy2026a{]} identifies the first link in the
chain: bounded active context \(C_n\) under information-preservation
pressure necessarily produces graph-structured knowledge. This paper
derives the remaining structural consequences as a connected chain from
the same generative constraint --- bounded \(C_n\) under survival
pressure in a world with positive entropy: the inseparability of
knowledge and action (Theorem 1), escalation dominance over top-down
allocation (Theorem 2a), evaluation-driven state revision (Theorem 3),
the encoding hierarchy (Theorem 4), type irreducibility (Theorem 5),
and the necessity of \(\mathcal{A}\) and \(\eta_M\) (Theorem 6). The
individual results formalize known intuitions; the contribution is that
they are all consequences of the same constraint, derived rather than
separately assumed.

\begin{center}\rule{0.5\linewidth}{0.5pt}\end{center}

\subsection{3. Formal Definitions}\label{formal-definitions}

\textbf{Definition 1 (World State Space).} Let \(W\) be a measurable
space of possible world states. At time \(t\), the true state
\(w_t \in W\) is not directly observable by any entity.

\textbf{Definition 2 (Knowledge Graph).} A knowledge graph \(K\) is a
directed graph where nodes are propositions about \(W\) and edges
represent inferential or causal relationships between propositions.
\(K\) is not a passive store; its retention criterion is established by
Theorem 1.

\textbf{Definition 2a (Bounded Knowledge Graph).} A knowledge graph
\(K\) has capacity \(K_{max}\): the maximum number of propositions \(K\)
can hold simultaneously, determined by the substrate's physical
constraints. When \(|K| < K_{max}\), \(K\) operates in the
\textbf{growth regime}: propositions accumulate freely subject to the
retention criterion of Theorem 1. When \(|K| = K_{max}\), \(K\) operates
in the \textbf{curation regime}: adding a new proposition requires
removing an existing one. The two regimes have different retention
criteria.

\textbf{Definition 3 (Informed Action Space).} Let
\(F \subseteq (K \times W_{obs} \to \Delta A)\) be a set of functions,
where \(W_{obs}\) is an entity's observable projection of \(W\) and
\(\Delta A\) is a distribution over possible actions. Each element
\(f \in F\) has type \(f : K \times W_{obs} \to \Delta A\). \(F\) is the
complete set of informed actions available to an entity. An
\textbf{informed action} \(f \in F\) is an action grounded in \(K\),
validated against \(W\), and attributed to its origins: one that knows
why it works. An uninformed action (grounded in no \(K\)) is not an
element of \(F\). Knowledge with no action potential is not an element
of \(F\) either; it is data. An informed action is specifically the
coupling: knowledge that enables action, action that is justified by
knowledge.

\(F\) carries graph structure. The elements of \(F\) are nodes; typed
relations between them are edges. This structure is not a design choice
--- it is forced by the same constraints (retrieval
\(O(|F^{[f]}|)\), bounded \(C_n\), deduplication) that force \(K\) to be
a directed graph (Theorem 2b Part 3), and independently by the
requirement for sub-linear retention evaluation with cascade (Theorem
2c). The proof is identical: any structure satisfying all constraints
simultaneously is a directed graph. Edge
types on \(F\) follow from the same evidence-grounded logic as edge
types on \(K\): hard dependencies between actions (feasibility
constraints), probabilistic ordering evidence (soft gradient
directions), and joint-execution synergies (combined actions that
outperform either alone).

\(F\) as a graph is discrete (nodes and typed edges); each individual
\(f \in F\) maps into the continuous probability simplex \(\Delta A\).
Gradient descent in \(F\) is traversal on the discrete graph, directed
by the continuous-valued \(G(f', K)\) at each adjacent node. \(K\) and
\(F\) are both directed graphs under bounded context and selection
pressure, differing only in domain: \(K\) is what the entity knows;
\(F\) is what it can do.

In general, \(K\) and \(F\) are distinct graphs with many-to-many
coupling. In the degenerate case where every proposition is directly
actionable, \(F = K\). This is the limit of crystallization at the
bottom of the encoding hierarchy: a reflex is simultaneously knowledge
and action. \(F = K\) holds locally in fully crystallized domains;
\(F \neq K\) persists globally because \(\mathcal{T}_{reachable}\) grows
without bound (Theorem 3, Part II).

\textbf{Definition 3a (Action Feedback).} For any \(f \in F\) executed
in world state \(w\), the feedback \(\phi(f, w) \in \Phi\) is the
observable consequence of \(f\) in \(w\): the signal \(W\) returns in
response to the action. The updated knowledge after executing \(f\) and
observing its feedback is:

\[K_f = K \cup \{p_f\}\]

where \(p_f\) is the proposition grounded by observing \(\phi(f, w)\).
The feedback \(\phi(f,w)\) is the formal object underlying the gradient
signal \(\delta\) in Theorem 3: \(\delta =
\phi(f, w) - \mathbb{E}[\phi(f, \cdot) \mid K]\).

\textbf{Definition 3b (Information Gain).} The information gain of
action \(f\) given current knowledge \(K\) is:

\[G(f, K) = \mathbb{E}_w\!\left[U(w, K) - U(w, K_f)\right] = H(W \mid K) - H(W \mid K,\, \phi(f, w)) = I(f\,;\, W \mid K)\]

\(G(f, K) \geq 0\) by the data processing inequality
{[}Cover2006{]}: conditioning on true information cannot increase
conditional entropy. \(G(f, K) = 0\) when \(f\) provides no new
information about \(w\) given \(K\). Corrupted feedback may increase
\(U\) and falls outside this definition's scope.

\textbf{Definition 4 (Entity).} An entity
\(E = (K, F, \sigma, \mathcal{A})\) where
\(\sigma : W \rightarrow W_{obs}\) is a sensing function projecting the
true world state into the entity's observable space. \(K\) and \(F\)
evolve over time via \(M\); their initial values \((K_0, F_0)\) are
inherited from \(L_0\) encoding. The bootstrapping problem is resolved
by \(L_0\): the entity begins from the encoded product of gradients that
preceded it. Throughout this paper \(K\) and \(F\) refer to
current-time values; the subscript \(0\) denotes initial conditions.

\textbf{Definition 5 (Uncertainty).} The uncertainty of entity \(E\)
about world state \(w\) is the surprisal of \(w\) given \(K\)
{[}Shannon1948, Cover2006{]}:

\[U(w, K) = -\log P(w \mid K)\]

\(K\) is a directed graph of propositions interpreted as a Bayesian
network {[}Pearl1988{]}. Each proposition \(p \in K\) contributes a
likelihood factor \(\ell_p : W \to [0,1]\). The graph's edge structure
encodes conditional independence: propositions with no path between them
are conditionally independent given shared ancestors. The posterior over
world states given the full knowledge graph is:

\[P(w \mid K) \propto P_0(w) \cdot \prod_{p \in K} \ell_p(w)\]

This is the standard Bayes posterior obtained by conditioning \(P_0\) on
all propositions simultaneously. \(P(w \mid K)\) is well-defined and
normalizable for any proper prior \(P_0\) and bounded likelihood
functions \(\ell_p\). The theory is parametric in \(P_0\) --- it does
not depend on a specific prior, only on its existence and
well-formedness.

The expected uncertainty over all world states is the conditional
entropy:

\[H(W \mid K) = \mathbb{E}_w[U(w, K)] = -\sum_w P(w \mid K) \log P(w \mid K)\]

The minimization imperative (Theorem 1) operates on this expected
quantity; \(U(w, K)\) is the pointwise uncertainty the entity faces when
it encounters specific world state \(w\).

\textbf{Definition 5a (Survival Threshold).} For entity \(E\) in domain
\(d \subseteq W\), there exists a threshold \(U_{lethal}^d\) such that:

\[U(w, K) > U_{lethal}^d \Rightarrow E \text{ does not survive in } d\]

Above this threshold, uncertainty is so high that the entity cannot act
effectively enough to persist. \(U_{lethal}^d\) is not a choice or a
preference; it is imposed by \(W\).

\textbf{Definition 6 (Agency).} Agency \(\mathcal{A} \in [0, 1]\) is the
probability that entity \(E\) engages its learning mechanism \(M\) in
response to a gradient signal from \(W\): the response gate. It is not
an external objective imposed on \(E\), nor a property derived from
\(F\) or \(K\). It is constitutive: \(\mathcal{A}\) is what makes \(F\)
an active informed action space rather than a static structure. At
\(\mathcal{A} = 0\), \(F\) is a set of possible functions with no
internal reason to engage \(W\). At \(\mathcal{A} = 1\), \(E\) engages
on every gradient signal it encounters. Intermediate values are
meaningful: the response probability can be higher or lower, shaped by
architecture, encoding history, and the density of unresolved
uncertainty in the entity's reachable \(W\). The informal gloss --- the
desire to know {[}Schmidhuber2010{]}, the intrinsic drive to reduce
uncertainty --- is a phenomenological description of this gate, not the
definition.

\(\mathcal{A}\) has an empirical operationalization:

\[\mathcal{A}(E) \propto \lim_{U_{lethal}^d \to \infty} \frac{dU}{dt}(E)\]

An entity with genuine \(\mathcal{A} > 0\) continues engaging \(M\)
when survival pressure is removed; an entity with \(\mathcal{A} = 0\)
stops. This isolates intrinsic engagement from survival-compelled
updating and is measurable in any system where survival pressure can be
varied independently of capacity.

\(\mathcal{A}\) precipitates from \(M\) by the same normalization
argument that precipitates \(K\) from \(F\). An entity that embeds its
directional drive in every application of \(M\) recomputes orientation
at every step --- the \(M\)-analog of the \(O(n \cdot k)\) scaling
failure. An entity that maintains \(\mathcal{A}\) as the shared
directional structure across all \(M\)-applications achieves \(O(1)\)
access to its drive. \(\mathcal{A}\) is therefore the first factor to
precipitate from \(M\)'s cross product: the normalized orientation that
every \(M\)-application shares, extracted for the same reason \(K\) is
extracted from \(F\).

\textbf{Definition 7 (Intelligence).} The intelligence of \(E\) is the
\emph{driven capacity to improve}: the exercise of \(\mathcal{A}\)
through \(M\) to refine \(F\) and \(K\) in response to gradient signals
{[}Schmidhuber2010, Thrun1998{]}:

\[I(E) = \mathcal{A} \cdot \eta_M\]

where \(\mathcal{A} \in [0, 1]\) is the intrinsic drive (Definition 6)
and \(\eta_M\) is the efficiency of the higher-order function space
\(M\) (Definition 11a): the rate at which \(M\) improves \(F\)'s
uncertainty-reduction capacity per unit of learning signal. This
expression is stated here as a definition and established as necessary
in Theorem 6: \(\mathcal{A}\) and \(\eta_M\) are the unique invariants
of persistence under bounded context and selection pressure, and their
product is the minimal measure that captures both.

At \(\mathcal{A} = 0\): \(I(E) = 0\) --- a capable but inert learning
mechanism is not intelligence. At \(\eta_M = 0\): \(I(E) = 0\) --- a
motivated entity whose \(M\) cannot improve \(F\) from evidence is an
automaton.

The expected uncertainty reduction along a trajectory is a
\textbf{performance measure} of current \(F\) and \(K\) --- the
consequence of past \(M\) activity, not the definition of intelligence.
Two entities starting from identical \(K\)
and \(F\) but different \(I(E)\) will diverge: the entity with higher
\(\eta_M\) builds better \(F\) faster, and eventually dominates in
\(U\)-reduction even if initially equal in performance.

Intelligence is not the size of \(K\) or \(F\). \(K\) and \(F\) are what
intelligence has built. Intelligence is the efficiency with which \(M\)
continues building.

\emph{Note on named components.} The components \(K\), \(F\), \(M\), and
\(\mathcal{A}\) named here are dimensions that have precipitated from
the entity's full cross-product state --- call it \(K^n\) --- at the
individual learning level. The theory's general result is that
dimensional precipitation must occur under the joint constraints: as
\(K^n\) grows, shared structure must be extracted and named or retrieval
cost grows without bound. What dimensions precipitate is a function of
the entity's \(W\) and its history; the theory constrains the process,
not the outcome. \(K\), \(F\), \(M\), and \(\mathcal{A}\) are the
dimensions that fall out in entities with active individual learning
mechanisms --- the class this paper studies. Other projections
precipitate differently: biological \(L_0\) encoding (DNA) does not
separate into \(K\), \(F\), \(M\), \(\mathcal{A}\); it is a dense
cross-product that carries all survival-relevant structure without
naming any of it as a distinct component. The named components are an
epistemological convenience grounded in the studied entity class, not a
universal constraint of the theory.

\begin{center}\rule{0.5\linewidth}{0.5pt}\end{center}

\subsection{4. The Inseparability of Knowledge and
Action}\label{the-inseparability-of-knowledge-and-action}

\textbf{Lemma 1 (Proposition Absence Lethality).} For any entity \(E\)
and proposition \(p \in K\) that grounds some \(f \in F\) in domain
\(d \subseteq W\): removing \(p\) from \(K\) weakly increases
uncertainty in \(d\), and in the limiting case where \(p\) is the sole
grounding for the critical action in \(d\), the increase may cause
\(U(w, K \setminus \{p\}) > U_{lethal}^d\), at which point \(E\) does
not survive in \(d\) (Definition 5a).

\emph{Proof.} Conditional entropy is monotonically non-increasing in the
conditioning set {[}Cover2006{]}:
\(H(W \mid K) \leq H(W \mid K \setminus \{p\})\) for any \(p \in K\).
Removing \(p\) weakly increases expected uncertainty. If \(p\) grounds
the only \(f \in F\) applicable in domain \(d\), then
\(K \setminus \{p\}\) contains no action grounded in \(d\): \(F\) has no
mechanism to reduce \(U(w, K \setminus \{p\})\) below its current level
in that domain when \(w \in d\) is encountered. By Definition 5a, if
\(U(w, K \setminus \{p\}) > U_{lethal}^d\) at the encountered \(w\),
\(E\) does not survive. Therefore the cost of absent \(p\) when needed,
\(C_{absent}\), is bounded below by entity non-survival in critical
domains --- which dominates all finite costs. \(\square\)

\textbf{Theorem 1 (K/F Inseparability).} For any entity
\(E = (K, F, \sigma)\), the retention criterion for any proposition
\(p\) is asymmetric:

\begin{itemize}
\tightlist
\item
  \(p\) is \textbf{retained} in \(K\) when \(\exists\, f \in F\)
  applicable to \(p\), or when the survival benefit of \(p\) is
  uncertain
\item
  \(p\) is \textbf{pruned} from \(K\) only when it is certain that \(p\)
  contributes zero to \(P(E, \tau)\) for any \(\tau\), i.e., when it is
  known that no \(f \in F\) can be grounded in \(p\), now or in any
  reachable region of \(W\)
\item
  \(p\) is \textbf{displaced} in the curation regime (\(|K| = K_{max}\))
  when a candidate proposition \(q\) satisfies
  \(\mathbb{E}[G(f_q, K)] > \mathbb{E}[G(f_p, K)]\), where \(p\) is the
  proposition with minimum marginal value in \(K\). Displacement
  requires only that \(q\) contributes more at the margin than \(p\) ---
  not that \(p\) is useless. The asymmetric retention criterion governs
  which propositions enter; the comparative criterion governs which are
  displaced when the graph is full. Information gain \(G\) is submodular
  --- the marginal value of \(p\) depends on what else is in \(K\), so
  the greedy displacement criterion is evaluated at the margin given the
  current full \(K\), not in isolation. This is the correct procedure:
  the entity estimates value relative to what it already holds, not
  relative to an empty graph {[}Nemhauser1978{]}.
\end{itemize}

The criterion is not symmetric. Uncertainty of benefit is not the same
as certainty of no benefit. An entity with \(\mathcal{A}\), whose
\(\mathcal{T}_{reachable}\) grows with every action taken, retains \(p\)
whose value it cannot yet see, because the gradient may reach that
territory. Pruning requires the same standard as certainty (Definition
10): the benefit must be known to be zero across sufficient variation in
\(W\).

\emph{Proof.} Let \(B(p, \tau)\) denote the contribution of \(p\) to
\(P(E, \tau)\) along trajectory \(\tau\), and let
\(\mathcal{T}_{reachable}\) be the set of trajectories reachable from
the entity's current state. Define the option value of retaining \(p\):

\[OV(p) = P\!\left(\exists\,\tau^* \in \mathcal{T}_{reachable} : B(p, \tau^*) > 0\right) \cdot \mathbb{E}\!\left[B(p,\tau^*) \mid B(p, \tau^*) > 0\right]\]

Let \(c_{retain} > 0\) denote the finite maintenance cost of holding
\(p\) in \(K\) (memory, retrieval overhead), and \(c_{prune} \geq 0\)
the cost of pruning \(p\) (processing, reorganization). Both are finite
by assumption: any physical knowledge store incurs bounded costs per
proposition. The expected net value of retention is
\(NV_{retain}(p) = OV(p) - c_{retain}\).

The expected net value of pruning is:

\[NV_{prune}(p) = -c_{prune} - P\!\left(\exists\,\tau^* : B(p,\tau^*) > 0\right) \cdot C_{absent}\]

where \(C_{absent}\) is the cost incurred when \(p\) is needed but
absent. By Lemma 1, \(C_{absent}\) is bounded below by entity
non-survival in critical domains, which dominates all finite costs:
\(C_{absent} \gg c_{retain}\).

Retention is selected for when \(NV_{retain}(p) > NV_{prune}(p)\). Since
\(C_{absent} \gg
c_{retain}\), this reduces to: retention holds whenever

\[P\!\left(\exists\,\tau^* \in \mathcal{T}_{reachable} : B(p,\tau^*) > 0\right) > 0\]

i.e., whenever survival benefit is uncertain. Even a small probability
of needing \(p\), multiplied by the cost of its absence, outweighs
finite maintenance.

Pruning holds only when \(P(\exists\,\tau^* : B(p,\tau^*) > 0) = 0\) ---
i.e., when it is certain that \(p\) contributes zero benefit across all
reachable trajectories. Any positive probability keeps the retention
condition active. The asymmetry is strict: uncertain benefit is
sufficient for retention; certain zero benefit is necessary for pruning.

In the curation regime (\(|K| = K_{max}\)), adding \(q\) requires
displacing \(p^* = \arg\min_{p \in K} \mathbb{E}[G(f_p, K)]\).
Displacement occurs when
\(\mathbb{E}[G(f_q, K)] > \mathbb{E}[G(f_{p^*}, K)]\): the asymmetric
criterion applied to a choice between two propositions rather than
between one and nothing. By Theorem 3 Part II,
\(\mathcal{T}_{reachable}\) grows monotonically, keeping future needs
uncertain unless benefit is certainly zero --- preserving the asymmetry
across the entity's lifetime. \(\square\)

\emph{Correspondence.} This formalizes Gibson's affordances
{[}Gibson1979{]} and extends the information bottleneck
{[}Tishby2011{]}, which treats the retention threshold as a design
parameter; here it is derived from survival pressure. The asymmetric
criterion resolves an apparent asymmetry:
an entity retains \(p\) beyond the certainty horizon because the
benefit is uncertain, not known to be zero.

\textbf{Corollary 1 (Indexed Knowledge).} \(K\) is constitutively
indexed by \(F\). Two entities with identical \(\sigma\) but different
\(F\) maintain different \(K\) even when perceiving identical world
states.

\textbf{Corollary 2 (Compactness).} \(K\) is bounded in size by \(F\):
only propositions applicable to available informed actions are retained.

\textbf{Corollary 2a (Curation and Forgetting).} Under bounded \(K\),
forgetting is correct behavior: the curation regime displaces
lower-value propositions as \(\mathcal{T}_{reachable}\) evolves,
ensuring \(K\) tracks the current frontier.

\textbf{Corollary 2b (K as the Normalization of F --- the Scaling
Argument).} An entity can maintain informed actions without a separately
structured \(K\), embedding all propositional knowledge implicitly in
the parameters of each \(f \in F\). This is a valid but scaling-limited
architecture. Let \(n = |F|\) and \(k\) = average number of shared
propositions per function. Without \(K\): context per action is \(O(k)\)
(each \(f\) carries its full background); update cost when one
proposition changes is \(O(n)\) (every \(f\) encoding it must be
individually revised); total cost scales as \(O(n \cdot k)\). With \(K\)
normalized out: context per action is \(O(|K^{[f]}|)\) --- the relevant
subgraph retrieved by traversal (Theorem 2b); update cost per
proposition is \(O(1)\) --- one update to \(K\) propagates to all \(f\)
that reference it; total cost scales as \(O(n + k)\).

The difference \(O(n \cdot k)\) versus \(O(n + k)\) is not a constant
factor --- it is the difference between a system whose context cost
grows multiplicatively with shared knowledge complexity and one whose
context cost stays bounded under growth of \(F\). Under bounded \(C_n\),
the implicit-\(K\) architecture does not scale: every increase in
\(|F|\) increases required context per action proportionally to the
shared background knowledge it carries. \(K\) precipitates from \(F\) as
the normalization of shared propositions --- the factoring-out of common
structure across functions --- because this is the only architecture
that keeps \(C_n\) bounded as \(F\) grows. An entity without \(K\) is
not wrong; it is limited to the regime where \(n \cdot k\) fits in
\(C_n\). Selection pressure at scale removes such entities in favor of
those whose \(K\) has crystallized --- literally, the precipitation of
shared knowledge out of the action space into its own normalized graph.
This is the formal explanation of the \(O(N \times M) \to O(1)\)
empirical measurement described in Section 1: each entity's local graph
traversal is \(O(1)\) precisely because shared propositions have been
factored out of \(F\) into \(K\).

\begin{center}\rule{0.5\linewidth}{0.5pt}\end{center}

\subsection{5. The Escalation Principle}\label{the-escalation-principle}

\textbf{Definition 8 (Encoding Hierarchy).} The encoding hierarchy is a
continuous two-dimensional space \(\mathcal{E}(c, r)\) parameterized by
runtime cost \(c \geq 0\) and reversibility \(r \in [0, 1]\), where
increasing \(c\) corresponds to more expensive engagement and increasing
\(r\) to more easily revised encodings. Consistent with the continuity
principle: a continuous being moves through this space continuously;
there is no categorical jump between encoding modes, only a gradient of
cost and permanence {[}Baars1988, Dehaene2011, Kahneman2011{]}.

Boundaries exist where the qualitative character of encoding changes,
requiring information to be translated as it crosses. These boundaries
are \textbf{interfaces}: formal structures through which signals pass
between qualitatively distinct encoding regimes. The named levels
\(L_0, L_1, \ldots, L_n\) are the interfaces at the significant
boundaries of \(\mathcal{E}(c, r)\):

{\def\LTcaptype{none} % do not increment counter
\begin{longtable}[]{@{}
  >{\raggedright\arraybackslash}p{(\linewidth - 6\tabcolsep) * \real{0.2500}}
  >{\raggedright\arraybackslash}p{(\linewidth - 6\tabcolsep) * \real{0.2500}}
  >{\raggedright\arraybackslash}p{(\linewidth - 6\tabcolsep) * \real{0.2500}}
  >{\raggedright\arraybackslash}p{(\linewidth - 6\tabcolsep) * \real{0.2500}}@{}}
\toprule\noalign{}
\begin{minipage}[b]{\linewidth}\raggedright
Interface
\end{minipage} & \begin{minipage}[b]{\linewidth}\raggedright
Boundary crossed
\end{minipage} & \begin{minipage}[b]{\linewidth}\raggedright
Runtime Cost
\end{minipage} & \begin{minipage}[b]{\linewidth}\raggedright
Reversibility
\end{minipage} \\
\midrule\noalign{}
\endhead
\bottomrule\noalign{}
\endlastfoot
\(L_n\) & Active reasoning / novel uncertainty & \(O(\text{compute})\) &
Fully reversible \\
\(\vdots\) & Learned patterns, habits & Decreasing & Decreasing \\
\(L_1\) & Reflex / autonomic & \(\approx 0\) & Semi-permanent \\
\(L_0\) & Heritable encoding & \(0\) & Permanent within lineage \\
\end{longtable}
}

Escalation (Theorem 2) is the \textbf{upward interface protocol}: the
specification of what form a failure signal takes when crossing a
boundary from lower \(c\) to higher \(c\). Promotion (Theorem 4) is the
\textbf{downward interface protocol}: the condition under which an
action may cross a boundary from higher \(c\) to lower \(c\). The cost
\(C_i\) in Theorem 4 is a property of the interface, not of the
continuous space between interfaces: it is the cost of failure when an
encoding error has crossed the boundary and propagates at the lower
level's permanence.

The number of interfaces is substrate-determined --- how many
qualitative boundaries exist in the entity's physical
\(\mathcal{E}(c, r)\). The theory specifies the structural properties
every interface must have (monotone cost, monotone rigor threshold,
promotion criterion); it does not prescribe how many interfaces exist or
what they are called. The named levels in the table above are empirical
illustrations, not definitional. The continuous space between interfaces
admits arbitrary encodings; the interfaces themselves are discrete
because a boundary is a specific location, not a region.

\textbf{Theorem 2 (Escalation).} An informed action \(f \in F\) is
handled at the lowest encoding level \(L_i\) such that
\(U(w, K^{[f]}) \approx 0\) for the domain of \(w\) relevant to \(f\),
where \(K^{[f]} \subseteq K\) is the knowledge subset applicable to
\(f\). Escalation to \(L_{i+1}\) occurs if and only if \(L_i\) fails to
bound uncertainty.

\emph{Proof.} By Definition 8, the cost of engaging level \(L_i\)
exceeds the cost of engaging \(L_{i-1}\): the encoding hierarchy is
parameterized by increasing runtime cost \(c\), so \(C_{i+1} > C_i\) for
all \(i\). Selection pressure therefore favors handling \(f\) at the
lowest level sufficient to bound uncertainty: engaging \(L_{i+1}\) when
\(L_i\) suffices incurs unnecessary cost without reducing \(U\) further.
An entity that routinely over-escalates incurs higher encoding costs
without survival benefit and is selected against. Conversely, if \(L_i\)
fails to bound \(U(w, K^{[f]})\) --- that is, \(U(w, K^{[f]}) > 0\)
beyond acceptable residual --- then \(f\) is not handled at \(L_i\): the
failure propagates upward until some \(L_j\) (\(j > i\)) contains the
uncertainty or the entity is harmed (Definition 5a). The active context
window \(L_n\) is the level of last resort by construction: it is the
highest-cost, fully reversible level. \(\square\)

\textbf{Corollary 3 (Context Window as Frontier).} \(L_n\) is the
frontier of unbounded uncertainty, not the seat of intelligence. Its
engagement signals failure in lower-level encodings.

\textbf{Corollary 4 (Simultaneous Multi-Level Activity).} All levels
operate concurrently: \(L_n\) handles novel uncertainty while lower
levels handle their domains independently.

\subsubsection{Why the Graph: Progressive Disclosure and the Bounded
Context
Theorem}\label{why-the-graph-progressive-disclosure-and-the-bounded-context-theorem}

\textbf{Definition 8a (Active Context Capacity and Inference Quality).}
Let \(C_n\) denote the capacity of \(L_n\): the maximum number of
propositions simultaneously active in context. \(C_n\) is bounded above
by the substrate's physical constraints --- but the effective capacity
available for any given problem is also shaped by the structure of
\(K\). Graph-structured knowledge enables chunking: a single graph node
can represent a cluster of propositions that co-activate as a unit,
making multiple propositions available at the cost of one context slot.
Expertise increases effective \(C_n\) within a domain by compressing
frequently co-occurring propositions into retrievable chunks --- not by
expanding the substrate bound, but by reducing the number of slots
required per unit of relevant knowledge. The substrate sets the ceiling;
the structure of \(K\) determines how much of \(W\) that ceiling covers.
Let \(\text{ctx}(L_n, t)\) be the set of propositions loaded into
\(L_n\) at time \(t\), with \(|\text{ctx}(L_n, t)| \leq C_n\).

The \textbf{inference quality} of \(L_n\) for problem \(f\) is the
relevance density of its content:

\[\rho(L_n, f) = \frac{|K^{[f]}|}{|\text{ctx}(L_n, t)|}\]

\(\rho = 1\): \(L_n\) holds exactly what \(f\) requires --- maximum
signal, zero noise. \(\rho \to 0\): \(L_n\) is dominated by content
irrelevant to \(f\) --- inference degrades. Every proposition in
\(\text{ctx}(L_n)\) that is not in \(K^{[f]}\) is noise: it contributes
variance to \(L_n\)'s inference without contributing signal. Minimizing
that noise is equivalent to minimizing
\(|\text{ctx}(L_n)| - |K^{[f]}|\).

\textbf{Definition 8b (Delegation Overhead).} Under a top-down
allocation architecture, \(L_n\) directs which subproblems are handled
at lower levels. For each active delegation, \(L_n\) must hold: the
subproblem specification, the expected return signal, and sufficient
context to integrate the result when it arrives. Let \(\alpha > 0\)
denote the minimum context cost per active delegation --- the overhead
top-down allocation imposes on \(L_n\) per managed subproblem,
regardless of whether that subproblem is routine or novel.

\textbf{Theorem 2a (Top-Down Allocation Degrades Inference Quality).}
Under a top-down allocation architecture with \(k(t)\) concurrently
delegated subproblems, the inference quality available to \(L_n\) for
genuinely novel problems satisfies:

\[\rho(L_n, f_{\text{novel}}) = \frac{|K^{[f_{\text{novel}}]}|}{|K^{[f_{\text{novel}}]}| + k(t)\cdot\alpha}\]

As the entity's informed action space \(F\) grows (more capability, more
concurrent processes), \(k(t)\) grows proportionally and \(\rho \to 0\).
The only recovery is to grow \(C_n\): the context window must scale
linearly with capability to maintain fixed \(\rho\).

\emph{Proof.} \(L_n\) holds \(K^{[f_{\text{novel}}]}\) for the current
problem plus \(k(t)\cdot\alpha\) for delegation overhead. Since
\(\alpha > 0\) and \(k(t)\) grows with \(|F|\), the denominator is
unbounded while \(|K^{[f_{\text{novel}}]}|\) remains bounded (any
specific problem requires only its relevant subgraph). Therefore
\(\rho \to 0\) as \(|F| \to \infty\). When
\(k(t)\cdot\alpha \geq C_n - |K^{[f_{\text{novel}}]}|\), \(L_n\) lacks
capacity for genuine inference entirely. To maintain \(\rho\) at a fixed
level, \(C_n\) must grow as \(k(t)\) grows. \(\square\)

\emph{Note on attention mechanisms.} The \(O(k)\) delegation overhead
applies to any mechanism that evaluates all \(n\) stored propositions to
determine relevance, including full self-attention (\(O(n^2)\) per
query). Sparse and learned attention mechanisms reduce this cost by
learning fixed relevance patterns --- which is learned graph
construction. See Theorem 2c Part 6 for the full argument that
continuous-weight retrieval mechanisms are implementations of graph
structure, not alternatives to it.

\emph{Note.} This is self-undermining on the Free Energy Principle's own
terms {[}Friston2010, Friston2017{]}. FEP holds that intelligent systems
minimize free energy --- equivalently, minimize the variance they must
manage. But growing \(C_n\) to accommodate delegation overhead adds
variance: more propositions in \(L_n\) means more uncertain relevance to
any specific problem, not less. The orchestrator architecture prescribed
by FEP's precision allocation is inconsistent with FEP's own
minimization imperative: allocating attention from \(L_n\) increases the
variance that \(L_n\) must hold in order to allocate.

\emph{Biological note.} We do not claim that biological precision
allocation operates like a token context window, nor that the cortical
hierarchy is unidirectional. Top-down signals exist at every level:
higher levels necessarily convey signals to lower levels, and the
bidirectionality of the cortical hierarchy is well established. The
architectural claim is narrower: it concerns \emph{engagement
initiation}, not signal direction. In an escalation architecture, what
triggers higher-level involvement is bottom-up failure --- a lower
level's inability to bound uncertainty --- not a top-down assignment of
work. Top-down signals flow as a consequence of engagement; they do not
determine when engagement occurs. The cost of \emph{top-down assignment
as the primary control architecture} is \(O(k)\) in high-level capacity,
as the attentional and dual-task literature demonstrates --- each
concurrently supervised subprocess consumes a share of the supervisory
resource. Escalation avoids this cost by making higher-level engagement
demand-driven rather than supply-driven. The transition from active to
automatic processing across skill acquisition is the biological
measurement of this convergence: as \(f\) becomes automatic (assigned to
\(L_0\)), it stops consuming high-level capacity, freeing \(L_n\) for
novel problems. The formal argument is substrate-neutral; biology is the
validation test.

\textbf{Theorem 2b (Graph-Structured Knowledge Bounds Active Context).}
Under the escalation architecture (Theorem 2) with \(K\) structured as a
directed graph (Definition 2), as \(K\) grows: (1) no delegation
overhead at \(L_n\); (2) inference quality \(\rho = 1\) (conditional on
edge fidelity); (3) \(C_n\) need not grow as \(|F|\) or \(|K|\) grows.

\emph{Proof summary.} Parts 1--2 follow from escalation: \(L_n\) holds
only the escalated problem's relevant subgraph \(K^{[f]}\), not
delegation context. Part 3 is graph necessity: the escalation
architecture and bounded context jointly impose retrieval
(\(O(|K^{[f]}|)\)), context (\(\rho = 1\)), and deduplication
constraints on \(K\), and graph structure is the only representation
satisfying all three without information loss. The complete proof ---
including the information-preservation argument (factoring vs.\
discarding vs.\ lossy compression), the retention-dynamics argument
(sub-linear dependency queries require stored adjacency), and the
demonstration that attention mechanisms implement graph structure rather
than constituting alternatives to it --- is given in
{[}McCarthy2026a{]}. Part 4: by Corollary 6, proven informed actions
promote to lower encoding levels, keeping \(L_n\) occupied only by
genuinely novel uncertainty. \(\square\)

\textbf{Corollary 4a (Selection Pressure Produces Graph-Structured
Knowledge).} Graph structure is not merely more efficient than
alternatives --- it is the necessary consequence of satisfying the
retrieval, context, and memory constraints jointly
{[}McCarthy2026a{]}. Entities that grow \(C_n\) to manage delegation
overhead are selected against because they violate the context constraint
as \(|K|\) grows; entities using approximate retrieval are selected
against because they violate \(\rho = 1\); entities duplicating shared
propositions are selected against because maintenance cost grows with
\(|F|\).

\textbf{Corollary 4b (The Scaling Argument Against Context Window
Growth).} Growing \(C_n\) directly is self-defeating {[}Liu2024{]}.
By Theorem 2a, more context means more noise: \(\rho\) decreases as
\(C_n\) fills with content not relevant to the current problem.
By Definition 7, intelligence is the rate at which \(M\) improves
\((F, K)\), not the capacity of \(L_n\). The efficient path is not to
grow \(C_n\) but to grow the encoded knowledge at lower levels ---
filling the graph, not the context window.

\begin{center}\rule{0.5\linewidth}{0.5pt}\end{center}

\subsection{6. Gradient Descent on
Uncertainty}\label{gradient-descent-on-uncertainty}

\textbf{Theorem 3 (Agency Necessarily Produces Evaluation-Driven State
Revision).}

\emph{On the space of descent and dimensionality.} The gradient descent
in this theorem operates in \(F\) --- the space of informed actions ---
not in \(W\). \(W\) may be arbitrarily rough, discontinuous, or of
unknown dimensionality; no assumption on \(W\)'s structure is required.
The entity does not descend over \(W\); it updates \(f \in F\) in
response to signals received from \(W\).

The dimensionality of \(W\) or of \(M_{\text{res}}\) is not assumed. We
do not know whether the latent space is finite or infinite-dimensional.
The convergence proof does not require this knowledge. It operates on
the scalar \(U(w, K) = H(w \mid K)\), which is bounded below by zero
regardless of the dimensionality of \(W\) or \(F\). The proof shows
\(\mathbb{E}[U]\) decreases monotonically under \(M\)-updates and
converges by the monotone convergence theorem, without decomposing \(U\)
into independent dimensions.

\emph{Part I: Forced Signal, Response Gate.} Every executed action
\(f_i\) in world state \(w\) generates a feedback signal
\(\delta = \phi(f_i, w) - \mathbb{E}[\phi(f_i, \cdot) \mid K]\)
(Definition 3a): the signed deviation of observed outcome from
expectation. This signal arrives from \(W\) but is applied in \(F\): it
is a subgradient of \(U\) as a functional on \(F\) at the current
\(f_i\), evaluated at the encountered \(w\). The signal is generated by
execution, not by \(\mathcal{A}\). It arrives whether or not the entity
has \(\mathcal{A}\). The entity does not choose whether to receive this
signal. \(W\) imposes it.

\(\mathcal{A}\) is the response gate: the probability that \(M\) engages
on \(\delta\) when \(\delta\) arrives. An entity with
\(\mathcal{A} = 0\) receives every signal but \(M\) does not engage;
\((F, K)\) remains unchanged. An entity with \(\mathcal{A} > 0\) engages
\(M\) on the signal, updating \((F, K)\) in response. The encountered
world states are on-policy --- determined by the entity's trajectory,
not sampled i.i.d. from \(P(w \mid K)\) --- but this does not impede
convergence. The subgradient \(\delta\) is valid at every encountered
\(w\) regardless of how \(w\) was reached; the entity is updating a
function in \(F\), and each \(\delta\) is a legitimate descent direction
in that space. The burn tells you that this \(f_i\) was insufficient at
\(w\).

An entity with \(\mathcal{A} = 0\) does not update in response to
\(\delta\). \(U\) does not decrease. It eventually exceeds
\(U_{lethal}^d\) in some critical domain. Selection removes it.
Therefore any entity that persists must have \(\mathcal{A} > 0\): not
because \(\mathcal{A}\) generates the signal, but because
\(\mathcal{A}\) is what enables the entity to act on signals that \(W\)
imposes on every execution regardless. When \(M\) engages:

\[\Delta f_i = -\eta \cdot \delta\]

\emph{Convergence of U.} The convergence argument does not decompose
by dimension. \(U(w, K) = H(w \mid K) \geq 0\) (Definition 5) is a
scalar quantity bounded below by zero. We show it decreases in
expectation under \(M\)-updates and use this to establish convergence
directly.

\textbf{Step 1 (Each update reduces U in expectation).} By Definition
3b, every informed action \(f \in F\) applied at world state \(w\) yields
information gain \(G(f, K) = I(f\,;\,w \mid K) \geq 0\), with equality
only when \(f\) provides no new information about \(w\) given \(K\).
When \(M\) engages on the feedback signal \(\delta\) and updates
\((F, K) \to (F', K')\), the updated knowledge \(K'\) satisfies
\(H(w \mid K') \leq H(w \mid K)\) --- conditioning on additional
evidence cannot increase entropy (data processing inequality,
{[}Cover2006{]}). Therefore:

\[\mathbb{E}[U(w, K') \mid K] \leq U(w, K) - G(f, K)\]

Each \(M\)-update that engages on a signal with \(G > 0\) strictly
reduces expected uncertainty. Updates with \(G = 0\) leave \(U\)
unchanged.

\textbf{Step 2 (Convergence).} The sequence
\(\{\mathbb{E}[U(w, K(t))]\}_{t \geq 0}\) is non-increasing (by Step 1)
and bounded below by zero. By the monotone convergence theorem,
\(\mathbb{E}[U]\) converges. Since
\(\sum_t G(f_t, K(t)) \leq U(w, K(0)) < \infty\), the information gains
\(G(f_t, K(t)) \to 0\): the marginal improvement per update vanishes
asymptotically. This is convergence to a local minimum of \(U\) given
the entity's current \(F\) and \(\sigma\).

\textbf{Step 3 (Dimensionality is irrelevant to convergence).} The
argument above operates on the scalar \(U\), not on individual
dimensions of \(W\) or components of \(F\). The entity's
parameterization of \(F\) may grow as
\(\mathcal{T}_{reachable}\) expands and new propositions are grounded
(Part II). This growth does not invalidate convergence: each new
proposition \(p_{new}\) added to \(K\) can only reduce \(U\) (by the
data processing inequality), and the cross-grounding of \(p_{new}\) with
existing propositions can reduce \(U\) further by improving the
resolution of knowledge already in \(K\). Adding structure to \(K\) is
analogous to adding a column to a matrix and re-normalizing: the
existing entries can become better resolved, not worse. The coupling
between new and existing knowledge \emph{accelerates} convergence; it
cannot prevent it, because the Lyapunov bound \(U \geq 0\) holds
regardless of dimensionality.

\emph{Part II: Non-termination.} Let \(\mathcal{T}_{reachable}(K)\) be
the set of trajectories reachable from the entity's current state given
knowledge \(K\) (as defined in Theorem 1). We claim that for any finite
\(K(t)\), the frontier of \(\mathcal{T}_{reachable}\) contains world
states \(w\) where \(H(w \mid K) > 0\), provided \(W\) is sufficiently
rich that each executed action can ground a genuinely new proposition.
In a finite \(W\) where all propositions are exhaustible, the frontier
could in principle be covered; the claim applies in any \(W\) where the
entity's epistemic horizon has not been closed --- which is the intended
domain.

By Definition 3a, executing any \(f \in F\) produces feedback
\(\phi(f, w)\) that grounds a new proposition \(p_f\) not previously in
\(K\), yielding \(K_f = K \cup \{p_f\}\). This richer \(K_f\) grounds
new informed actions not available from \(K\) alone (Definitions 3 and
3b): actions that require \(p_f\) to select or apply, reaching world
states not reachable without \(p_f\). Therefore
\(\mathcal{T}_{reachable}(K_f) \supseteq
\mathcal{T}_{reachable}(K)\), strictly whenever \(p_f\) grounds at least
one new \(f' \in F\). Since Part I establishes that every executed
action generates a feedback signal, \(\mathcal{T}_{reachable}\) grows
monotonically with \(K\). At any finite time \(t\), \(K(t)\) has not
covered \(\mathcal{T}_{reachable}(K(t))\): the frontier always contains
states \(w\) where \(H(w \mid K) > 0\). Therefore \(\nabla_F U\) is
never globally zero over \(\mathcal{T}_{reachable}\), and descent never
terminates within the entity's reachable epistemic domain.

\emph{Proof.} Part I: \(U \geq 0\) is a Lyapunov function; each
\(M\)-update with \(G > 0\) strictly reduces \(\mathbb{E}[U]\) (data
processing inequality); \(\mathbb{E}[U]\) converges by monotone
convergence. Entities with \(\mathcal{A} = 0\) do not engage \(M\) and
are removed by selection. Part II: follows from the monotonic growth of
\(\mathcal{T}_{reachable}\) with \(K\) --- each executed action extends
the reachable frontier, ensuring new unexplored states are always
accessible within the domain of the entity's possible experience.
\(\square\)

\textbf{Corollary 5 (Evaluation-Driven State Revision is Survival, Not
Metaphor).} For individual learning, gradient descent on \(U\) is not a
description of how learning happens to work. It is what survival-driven
updating \emph{is}: the entity that burns and does not become more
careful does not survive. The entity that survives is, by definition,
the one whose \(M\) engaged on the signal. At the evolutionary
timescale, the mechanism differs in kind: random mutation perturbs the
\(L_0\) encoding, generating new \((F, K, M, \mathcal{A})\)
configurations; \(W\) evaluates each through the survival criterion;
configurations that persist reproduce. No individual organism follows a
gradient --- the gradient is revealed by differential survival of
randomly generated candidates. This is structurally the same as a
locally irrational action (temporarily negative \(\eta_M\)) at the
individual level: a perturbation that generates a new starting point for
the \(M \times \mathcal{A} \times K\) cross-product to be evaluated by
\(W\). The two mechanisms occupy a spectrum: at one end, random mutation
with binary survival feedback (evolution); at the other, directed update
via continuous \(\delta\) carrying full gradient information (individual
learning). Gradient descent is the efficient limit of this spectrum
under two conditions: (i) \(\delta\) carries directional information in
\(F\) --- not merely pass/fail, but a signed deviation identifying which
aspect of \(f_i\) was insufficient at \(w\); and (ii) \(M\) can use that
directional information to bias the update toward lower \(U\) rather
than generating a random candidate. When both conditions hold, each
evaluation step directly reduces \(U\). When neither holds, the
mechanism is random search with \(W\)-evaluation. Real learning systems
occupy positions between these limits; signal informativeness determines
how closely the process approximates directed gradient descent.

\textbf{Remark (Multi-timescale Structure of Evaluation-Driven
Revision).} Individual learning and evolutionary selection share a
common structure: execute, receive signal, retain or revise. The
mechanisms and mathematical objects differ in kind.

For individual learning, the mechanism is directed: the gradient signal
\(\delta\) arrives from a specific failure, \(M\) updates \((F, K)\) in
the direction that reduces \(U\):

\[\Delta F_{\text{individual}} = -\eta_{\text{ind}} \nabla_F \, U(w, K)\]

For evolutionary selection, the mechanism is undirected: random mutation
perturbs the \(L_0\) encoding of \((F, K, M, \mathcal{A})\), generating
candidate configurations. \(W\) evaluates each through survival. The
aggregate population-level distribution shifts toward configurations
with lower \(\mathbb{E}[U]\), but no individual organism descends a
gradient --- the gradient is revealed by differential survival of
randomly generated candidates. The notation
\(\Delta F_{\text{species}}\) denotes the shift in the population
distribution over \(F\)-variants, not a gradient step by any individual:

\[\Delta F_{\text{species}} = \operatorname{select}(\operatorname{mutate}(F_{\text{population}}))\]

Both processes are instances of the same structure: generate a candidate
state, submit to \(W\)'s evaluation, retain if valid. Gradient descent
(\(\Delta F_{\text{individual}}\)) is the efficient limit when signals
are informative enough to direct updates. Mutation followed by selection
(\(\Delta F_{\text{species}}\)) is the same mechanism without the
directed signal. This structural equivalence is not an identity ---
mechanisms, timescales, and mathematical objects are all distinct ---
but it grounds the common observation that both processes move toward
lower \(U\) at their respective timescales, and grounds a provable
consequence in Part II.

\emph{Derivation.}

\textbf{Part I (Common structure, different mechanisms).} Individual
learning updates \(F\) for a single entity based on \(U(w, K)\) at
specific world states \(w\) encountered by that entity. Evolutionary
change operates on a different object: it updates the population-level
distribution over \(F\)-variants through differential reproduction of
randomly generated candidates. These are not the same mathematical
operation --- individual learning changes a function within one agent
via directed update; evolutionary change changes a probability
distribution over many agents via undirected mutation and selection. The
common structure is: generate a new internal state, submit it to \(W\)'s
evaluation, retain if valid. Gradient descent is the efficient limit of
this structure; random mutation followed by selection is the general
case.

\textbf{Part II (\(\eta_{\text{evo}} \ll \eta_{\text{ind}}\) follows
from Theorem 4).} Although the mechanisms differ, both rates can be
characterized by the rigor threshold of their respective encoding
levels. The learning rate \(\eta_i\) at encoding level \(L_i\) is
inversely proportional to the rigor threshold \(\theta_i\): a higher
threshold requires more evidence per unit update, which is the
definition of a lower update rate. Formally:

\[\eta_i \propto \frac{1}{\theta_i} = \frac{1}{1 - \varepsilon/C_i} = \frac{C_i}{C_i - \varepsilon}\]

For individual learning at \(L_n\): \(C_n\) is the cost of a single
wrong action, relatively small, so \(\varepsilon/C_n\) is
non-negligible, \(\theta_n = 1 - \varepsilon/C_n\) is meaningfully less
than 1, and \(\eta_{\text{ind}} = C_n/(C_n - \varepsilon)\) is
correspondingly large. For evolutionary learning at \(L_0\): \(C_0\) is
lineage survival, far exceeding \(\varepsilon\), so
\(\theta_0 = 1 - \varepsilon/C_0 \approx 1\) and
\(\eta_{\text{evo}} \approx 1\). Since
\(\theta_0 \approx 1 \gg \theta_n\), it follows that
\(1/\theta_0 \ll 1/\theta_n\), therefore
\(\eta_{\text{evo}} \ll \eta_{\text{ind}}\). The learning rate
difference is not assumed; it is derived from the cost-of-failure
difference established in Theorem 4.

An individual may establish an informed action that reduces \(U\) in one
context but fails in others. Only informed actions that reduce \(U\)
across sufficient variation in \(W\) survive the evolutionary test ---
the higher rigor threshold filters for robustness. The analogy captures
this shared directional property: both processes favor what reduces
\(U\) at the relevant scale. The timescale, the rigor threshold, the
object being updated, and the mechanism all differ; the directional
relationship to \(U\) is what the analogy preserves.

\textbf{Theorem 4 (Encoding Permanence Scales with Rigor).} The
certainty threshold required to promote an informed action \(f\) to
encoding level \(L_i\) is:

\[\theta_i = 1 - \frac{\varepsilon_{\text{acceptable}}}{C_i}\]

where \(C_i\) is the cost of failure at level \(L_i\) and
\(\varepsilon_{\text{acceptable}}\) is the maximum acceptable expected
cost of an encoding error, expressed in the same units as \(C_i\). For
heritable encoding (\(L_0\)), \(C_0\) is lineage survival, requiring
near-certainty across the full distribution of environments. For active
reasoning (\(L_n\)), \(C_n\) is the cost of a single wrong action, a far
lower bar.

\emph{Proof.} Encoding \(f\) at level \(L_i\) is irreversible at cost
\(r_i > 0\) (Definition 8). The expected cost of encoding an action with
current confidence \(\theta\) is:

\[\mathbb{E}[\text{encoding error cost}] = (1 - \theta) \cdot C_i\]

where \((1-\theta)\) is the probability the action is wrong and \(C_i\)
is the cost when it is. Promotion is safe when this expected cost does
not exceed what is acceptable:

\[(1 - \theta) \cdot C_i \leq \varepsilon_{\text{acceptable}}\]

Rearranging:

\[\theta \geq 1 - \frac{\varepsilon_{\text{acceptable}}}{C_i}\]

The minimum threshold at which promotion is safe is therefore
\(\theta_i = 1 - \varepsilon_{\text{acceptable}}/C_i\). Below this
threshold, the expected cost of encoding error exceeds
\(\varepsilon_{\text{acceptable}}\) and continued active engagement is
preferred. At or above it, encoding is the lower-expected-cost choice.
\(\square\)

\emph{Consequence.} You must be very certain to permanently change the
source code. A bad encoding at \(L_0\) does not kill one individual; it
kills the lineage. The rigor of the test is proportional to the
permanence of the consequence.

\emph{Note on error propagation through the hierarchy.} A potential
objection is that the proof derives \(\theta_i\) for each level
independently, whereas errors in a hierarchy compound: a failure at
\(L_0\) propagates upward. This objection assumes a naive cascading
model in which errors at \(L_i\) arrive at \(L_{i+1}\) at rate
\((1-\theta_i)\) regardless of whether \(L_i\) could address them. That
is not the escalation architecture.

Under escalation (Theorem 2), errors partition into two classes at each
level: \textbf{handled errors} --- failures \(L_i\) detects and resolves
within its scope, which never reach \(L_{i+1}\) --- and
\textbf{unhandled exceptions} --- failures outside \(L_i\)'s domain or
beyond its capacity to resolve, which escalate. The rate seen by
\(L_{i+1}\) is not the compounded Bernoulli rate from below. It is the
rate of genuinely irreducible errors at \(L_i\): errors that \(L_i\) was
correct not to handle, because the knowledge or capacity to resolve them
does not exist at \(L_i\). This is a different quantity from
\((1-\theta_i)\), and it is bounded by the design of \(\theta_i\) itself
--- the threshold ensures \(L_i\) handles everything within its
cost-justified scope.

Furthermore, when an unhandled exception escalates, the context at
\(L_i\) was already disrupted by the failure: the action failed, and
\(L_i\)'s current trajectory is suspended or invalid. The attention cost
at \(L_{i+1}\) is not paid on top of productive work; it is paid instead
of broken work. Escalation is correct triage. The thresholds
\(\theta_i\) are therefore independent by construction: each level's
\(\varepsilon_{\text{acceptable}}\) is its own irreducible error domain,
not an accumulation from below. The proof applies at each level without
modification.

\textbf{Definition 10 (Certain).} An informed action \(f\) is
\textbf{certain} when \(U(w, K_f) \approx 0\) across sufficient
variation in \(W\) to meet the promotion threshold \(\theta_i\) (Theorem
4). Certainty is asymptotic: decreasing density of uncertainty until
active engagement no longer justifies the residual gradient.

\textbf{Corollary 6 (Bounded Adaptation).} Certain informed actions
escape the active context window permanently, freeing \(L_n\) for the
next frontier of uncertainty.

\begin{center}\rule{0.5\linewidth}{0.5pt}\end{center}

\subsection{7. The Higher-Order Function
Space}\label{the-higher-order-function-space}

\textbf{Definition 11 (Higher-Order Function Space).} Let \(M\) be a
function space operating on the joint \((F, K)\) pair {[}Thrun1998{]}:

\[M : (F, K) \times \mathcal{L} \rightarrow (F, K)\]

where \(\mathcal{L}\) is a learning signal (experience, execution
results, selection pressure). \(M\) is the mechanism by which
\((F, K)\) changes --- categorically different from \(F\), not a more
certain version of it.

\(M\)'s type signature is complete without enumerating its internal
structure, because \(M\) is the cross product of all dimensions that
have not yet precipitated out at the current scale of the entity:

\[M = \prod_{i \in \text{remaining}} D_i\]

where the product ranges over all dimensions --- drive, meta-learning
rate, exploration policy, convergence criteria, and whatever else has
not yet reached the scale threshold to become a named dimension. \(F\),
\(K\), and \(\mathcal{A}\) are factors that have already been extracted
from this product. \(M\) is always the remainder: the undifferentiated
joint space of everything not yet named. As the entity scales,
additional dimensions precipitate out of \(M\) by the same normalization
argument that precipitated \(K\) from \(F\): when the shared structure
of a dimension across \(M\)-instances is large enough that factoring it
out reduces retrieval cost from \(O(n \cdot k)\) to \(O(n + k)\),
selection pressure forces the extraction. Each precipitation shrinks
\(M\) by one factor and names a new dimension.

The hierarchy is open upward: there is always potentially a \(M^{(2)}\)
that operates on \(M^{(1)}\) the same way \(M^{(1)}\) operates on
\((F, K)\). One level suffices until it does not --- when the gradient
can no longer be followed by improving \((F, K)\) alone, pressure
propagates upward. The required number of levels is determined by the
problem, not prescribed by the framework. The precipitation argument
applies at every level; the hierarchy is not a prescribed structure but
the output of a single recursive process applied at increasing scale.
See Open Question 7.

The degenerate cases of level merger generalize from the \(F = K\) case:
for any contiguous block of \(N\) precipitated dimensions together with
the \(M\) residual, all \(N+1\) collapse to a single undifferentiated
process when the domain is fully crystallized at all scales in that
block. The merger condition is \(U(w, K) \approx 0\) across the domain
at all \(N\) scales simultaneously --- one condition, not \(N\) pairwise
conditions. Since \(M\) represents the full cross product of remaining
dimensions, any merger including \(M\) collapses not just \(N\) named
levels but the entire unprecipitated hierarchy simultaneously. \(F = M\)
is the maximally primitive state: action spans the full cross product of
all remaining dimensions. \(M = M^{(2)}\) means the cross product is
self-similar under the meta-operation --- evolution approaches this
limit at the timescale of selection on selection mechanisms.

\textbf{Definition 11a (M Efficiency).} The efficiency \(\eta_M\) of
\(M\) is the rate at which \(M\) improves the joint \((F, K)\)
capability per unit of learning signal:

\[\eta_M(E) = \mathbb{E}_{\mathcal{L}} \left[ \frac{\max_{f' \in F'} G(f', K') - \max_{f \in F} G(f, K)}{|\mathcal{L}|} \right]\]

where \((F', K') = M((F, K), \mathcal{L})\) is the updated joint state
after \(M\) acts, \(G(f, K) = I(f\,;\, W \mid K)\) is the information
gain of the best available informed action (Definition 3b) in bits, and
\(|\mathcal{L}|\) is the information content of the learning signal in
bits: \(|\mathcal{L}| = H(\mathcal{L})\), the entropy of the signal
received. Both numerator and denominator are in bits; \(\eta_M\) is
dimensionless, measuring improvement in best-available information gain
per bit of learning signal. \(\eta_M\) is defined for
\(H(\mathcal{L}) > 0\). When \(H(\mathcal{L}) = 0\), the learning signal
is degenerate (constant; carries no information) and \(\eta_M\) is
undefined: the entity received no learnable signal and the rate of
improvement per bit is not defined.

\(\eta_M\) is a signed quantity ranging over \((-\infty, \infty)\).
\(\eta_M > 0\): \(M\) improves the best available action in expectation
--- the entity learns. \(\eta_M = 0\): \(M\) cannot improve \((F, K)\)
on average --- the entity executes without learning. \(\eta_M < 0\):
\(M\) degrades \((F, K)\) locally --- overfitting, catastrophic
forgetting, destructive updating, or deliberate exploration. An entity
with \emph{persistently} \(\eta_M < 0\) does not persist under selection
pressure. An entity that can act on \(\eta_M < 0\) \emph{temporarily}
--- when local gradient following would trap it --- has strictly higher
long-run persistence probability than one constrained to
\(\eta_M \geq 0\) at every step. \(\eta_M\) is operationally independent
of \(\mathcal{A}\): one measures the mechanism of improvement; the other
measures the drive to engage it.

\emph{On measurement.} \(\eta_M\) as defined requires computing
\(G(f, K) = I(f\,;\, W \mid K)\), which depends on the true joint
distribution \(P(W, \phi(f,w))\) --- the distribution the entity is
trying to learn. The entity cannot compute its own \(\eta_M\) exactly.
This is not a defect of the definition but a structural feature:
\(\eta_M\) is a property of the entity-world coupling, not a quantity
the entity accesses directly. The same is true of \(\mathcal{A}\), whose
operationalization (Definition 6) requires varying survival pressure
externally. Both quantities are measurable by an observer with access to
the entity's trajectory and outcomes --- the standard setting for
evaluating learning systems --- but not by the entity itself in real
time. The theory's predictions (Theorems 1, 2, 4, 6) do not require
the entity to compute \(I(E)\); they follow from the structural
constraints regardless of whether the entity knows its own intelligence.

\emph{Remark on exploration and the necessity of acting on negative
\(I(E)\).} Theorem 3 Part II establishes that
\(\mathcal{T}_{reachable}\) grows without bound: the gradient landscape
is non-stationary and non-convex. An entity constrained to
\(\eta_M \geq 0\) at every step is a pure gradient follower --- it
converges to the nearest local optimum and remains there. But local
optima become maladaptive as new uncertainty domains emerge: an entity
trapped at a local optimum will eventually encounter
\(U > U_{lethal}^d\) in a domain it cannot reach from where it is.
Selection removes it.

Therefore temporarily negative \(\eta_M\) --- locally irrational steps
to escape local optima --- is survival-optimal when persistence requires
\(\mathbb{E}_t[\eta_M] > 0\) over the trajectory, not
\(\eta_M \geq 0\) at every step. Exploration and creative destruction
are the signature of an \(M\) sophisticated enough to optimize over
trajectories rather than steps.

\(M\) is itself subject to gradient descent at the meta-learning
timescale, with the same rigor thresholds (Theorem 4). \(F\) acts on
\(W\); \(M\) acts on \((F, K)\).

\textbf{Theorem 5 (Type Irreducibility).} No element of \(F\) can become
an element of \(M\) through any learning process, regardless of
certainty achieved.

\emph{Proof.} \(F : K \times W_{obs} \rightarrow \Delta A\).
\(M : (F, K) \times \mathcal{L} \rightarrow (F, K)\). These functions
have different type signatures and operate over different domains. No
learning process changes the type of a function. An informed action
\(f \in F\) that achieves certainty is promoted to a lower encoding
level; it does not change type. \(M\) takes \((F, K)\) as input and
returns an updated \((F, K)\); \(f\) takes knowledge and observations as
input and returns an action distribution. These are irreducibly
distinct. \(\square\)

\textbf{Corollary 5a (Distinct Roles).} \(F\) navigates \(W\) given
\(K\). \(M\) evolves \((F, K)\). Neuroplasticity is \(M\) acting on
\(F\); insight is \(M\) acting on \(K\); habit formation is \(M\)
promoting from \(L_n\) to \(L_2\); evolution is \(M\) acting on \(M\)
itself at the highest rigor threshold.

\textbf{Theorem 6 (Necessary Conditions for Persistence: \(\mathcal{A}\)
and \(\eta_M\) as Structural Invariants).} For any entity \(E\)
persisting in \(W\) under selection pressure \(U_{lethal}^d\) with
bounded context \(C_n\), in a world with positive entropy (one that
generates novel configurations over time):

\begin{enumerate}
\def\labelenumi{\arabic{enumi}.}
\tightlist
\item
  \(\mathcal{A}(E) > 0\) necessarily
\item
  \(\eta_M(E) > 0\) necessarily
\item
  Every other scalar property of \(E\) can equal zero for some
  persistent entity
\item
  Therefore \(I(E) = \mathcal{A} \cdot \eta_M\) is the natural measure
  of persistence --- identified by what survival requires, not
  stipulated
\end{enumerate}

\emph{Scope note.} This theorem establishes that \(\mathcal{A} > 0\)
and \(\eta_M > 0\) are \emph{necessary conditions} for persistence. The
proof is a selection argument: it shows that entities lacking either
property are eliminated. This explains why \(\mathcal{A} > 0\) entities
predominate --- it does not explain the \emph{origin} of
\(\mathcal{A}\), i.e., how self-maintaining structures with
\(\mathcal{A} > 0\) bootstrap from initial conditions where no such
structures exist. The origin question is a distinct problem, addressed
as an open conjecture in Section 11 (Open Question 1). The framework
takes persistence as its starting analytical condition and derives
structural requirements from it; it does not claim to explain why
anything persists in the first place.

\emph{Proof.}

\textbf{Part 1 (Necessity of \(\mathcal{A} > 0\)).} Suppose
\(\mathcal{A}(E) = 0\). By Definition 6, \(E\) does not respond to
gradient signals from \(W\): \((F(t), K(t)) = (F(0), K(0))\) for all
\(t\). \(W\) is a physical state space with positive entropy: it
generates novel configurations that no finite \(K(0)\) anticipates. This
is not a modeling assumption but a thermodynamic fact --- any \(W\) with
energy gradients produces states not represented in a fixed, finite
knowledge graph. Since \(K(0)\) is fixed and was not constructed with
these novel states in mind, \(U(w, K(0))\) is unconstrained there.

The selection argument does not require that any single entity
encounters lethal novelty within its lifetime. It operates on lineages.
An entity with \(\mathcal{A} = 0\) reproduces copies of itself ---
copies that inherit the same fixed \((F(0), K(0))\). Those copies
disperse across \(W\). In a world with positive entropy, the lineage's
collective coverage of \(W\) grows with each generation, even though
each individual's \((F, K)\) is static. The lineage therefore encounters
the novel configurations that \(W\) generates: states where
\(U(w, K(0)) > U_{lethal}^d\) (Definition 5a). No copy can adapt
(\(\mathcal{A} = 0\)); copies in those regions are removed. Copies in
stable niches persist temporarily, but the niches themselves are not
permanent --- \(W\) has positive entropy, so no region stays static
indefinitely. The \(\mathcal{A} = 0\) lineage contracts as its
environment shifts until it is eliminated. Selection removes the lineage.
Any lineage that persists must contain entities that update \((F, K)\)
--- i.e., entities with \(\mathcal{A} > 0\).

The mechanism by which \(\mathcal{A} > 0\) is enforced is
timescale-specific. At the individual learning timescale: entities with
\(\mathcal{A} = 0\) do not engage \(M\) on gradient signals and are
directly removed. At the evolutionary timescale: lineages that fail to
produce offspring with \(\mathcal{A} > 0\) are eliminated through
differential reproduction --- no individual requires \(\mathcal{A}\)
explicitly; the population-level filter enforces the property on
lineages. At the cultural timescale: groups that do not maintain the
institutional equivalent of \(\mathcal{A}\) (the drive to reduce shared
uncertainty) are outcompeted. The necessity is universal; the
enforcement mechanism is timescale-specific. \(\square\)

\textbf{Part 2 (Necessity of \(\eta_M > 0\)).} Suppose
\(\mathcal{A}(E) > 0\) but \(\eta_M(E) \leq 0\).

\emph{Case \(\eta_M = 0\):} \(E\) has the drive to reduce \(U\) but
\(M\) produces no net improvement in \((F, K)\) per unit of learning
signal: the best available action's information gain does not increase
after \(M\) acts. \((F, K)\) is effectively fixed despite
\(\mathcal{A}\). This reduces to Part 1's condition: with \((F, K)\)
fixed, the frontier of \(\mathcal{T}_{reachable}\) grows beyond \(K\)'s
coverage. Additionally, \(C_n\) is bounded (Definition 8a). As \(E\)
encounters novel world states, \(L_n\) accumulates unresolved context.
By Theorem 2a, inference quality \(\rho\) degrades as \(L_n\) fills
without promotion of resolved actions to lower encoding levels.
Promotion requires \(\eta_M > 0\): Theorem 4 establishes that encoding
at level \(L_i\) requires rigor threshold
\(\theta_i = 1 - \varepsilon/C_i\), and achieving that threshold
requires \(M\) to improve the action to the required certainty. With
\(\eta_M = 0\), no action ever reaches its threshold: \(L_n\) is never
relieved, \(\rho \to 0\), and inference eventually fails. \(U\) exceeds
\(U_{lethal}^d\) in the domains that required the new actions. Selection
removes \(E\).

\emph{Case \(\eta_M < 0\) persistently:} \(M\) degrades \((F, K)\) on
average across the trajectory. This accelerates the failure of the
\(\eta_M = 0\) case: \((F, K)\) worsens in expectation, \(U\) approaches
\(U_{lethal}^d\) from above. Selection removes \(E\) strictly faster.
Note that temporary \(\eta_M < 0\) --- deliberate exploration to escape
local optima --- is compatible with persistence and can increase
long-run survival probability (see Remark in Definition 11a). The
persistence condition is \(\mathbb{E}_t[\eta_M] > 0\) over the
trajectory, not \(\eta_M(t) > 0\) at every step.

Therefore any \(\mathcal{A} > 0\) entity that persists must have
\(\mathbb{E}_t[\eta_M] > 0\). \(\square\)

\textbf{Part 3 (Minimality: no other property is invariant).} For every
other scalar property of \(E\), there exists a persistent entity for
which it equals zero or its minimum value:

\begin{itemize}
\tightlist
\item
  \emph{\(K(0) = \emptyset\):} An entity with no initial knowledge
  survives if \(\mathcal{A} > 0\) and \(\eta_M > 0\) --- it builds \(K\)
  from gradient descent, beginning with the first feedback signal
  received from \(W\) (Theorem 3, Part I).
\item
  \emph{\(F(0) = \{\text{id}\}\):} \(F\) is a function space; its
  minimal element is the identity --- the action that changes nothing.
  The identity element exists independently of \(K\): it requires no
  grounding because it makes no knowledge claim. Executing the identity
  against \(W\) still returns feedback --- the world observed in the
  absence of action is itself a signal. That feedback grounds the first
  proposition \(p \in K\), after which the first non-trivial informed
  action can be constructed. The minimally non-empty \(F\) is therefore
  \(\{\text{id}\}\), not \(\emptyset\), and it is sufficient to begin
  gradient descent.
\item
  \emph{\(C_n = C_{min}\):} An entity with the smallest viable context
  window survives if the graph structure (Theorem 2b) keeps \(\rho = 1\)
  --- \(C_n\)'s absolute size is not invariant, only its boundedness.
\item
  \emph{Substrate, sensing \(\sigma\), world projection \(W_{obs}\):}
  These are free parameters of Definition 4 and Definition 1. The theory
  is parametric in all of them; no specific value is required for
  persistence. Entities on radically different substrates (biological,
  computational, collective) can all satisfy the survival constraint.
\end{itemize}

No scalar property other than \(\mathcal{A}\) and \(\eta_M\) has the
survival-critical necessity shown in Parts 1 and 2. \(\square\)

\textbf{Corollary 6a (\(I(E) = \mathcal{A} \cdot \eta_M\) is
structurally necessary, not stipulated).} From Theorem 6,
\(\mathcal{A}\) and \(\eta_M\) are the unique invariants of persistence
under bounded context and selection pressure. Any measure of
intelligence must be positive if and only if both are positive, and zero
if either is zero.

The product is not a choice --- it is the only normalized combination of
a probability and a conditional rate. \(\mathcal{A} \in [0,1]\) is the
probability that \(E\) engages \(M\) in response to a gradient signal.
\(\eta_M\) is the rate of improvement conditional on \(M\) being
engaged. By the law of total expectation:

\[\mathbb{E}[\text{improvement rate}] = \mathcal{A} \cdot \eta_M + (1 - \mathcal{A}) \cdot 0 = \mathcal{A} \cdot \eta_M\]

Any other combination fails: \(\mathcal{A} + \eta_M\) assigns positive
intelligence to a zero-drive entity (contradicting Theorem 6 Part 1);
\(\min(\mathcal{A}, \eta_M)\) erases the independent invariance of Parts
1 and 2. The product is the same structure as Theorem 1's retention
criterion: probability of needing it times conditional value.

\begin{center}\rule{0.5\linewidth}{0.5pt}\end{center}

\subsection{8. The Unified Model}\label{the-unified-model}

\textbf{Definition 12 (Trajectory).} A trajectory \(\tau\) is a sequence
of informed actions drawn from \(F\):

\[\tau = (f_1, f_2, \ldots, f_n), \quad f_i \in F\]

The \textbf{performance} realized by following trajectory \(\tau\) is:

\[P(E, \tau) = \mathbb{E}_W\left[U(w_0, K_0) - U(w_n, K_n)\right]\]

This is the total uncertainty reduction achieved --- a measure of what
current \((F, K)\) accomplished, not of intelligence as defined.
\(P(E, \tau)\) grows with trajectory length and with the quality of
\(F\) at the time of execution. Intelligence \(I(E) = \mathcal{A}
\cdot \eta_M\) is what determines how quickly \(P\) improves across
successive trajectories: an entity with high \(I(E)\) builds better
\((F, K)\) between trajectories, so \(P\) grows faster over time than
for an entity with low \(I(E)\) starting from identical initial
conditions.

Agency manifests as a trajectory \(\tau\): a sequence of informed
actions selected and executed over time. Each action reduces \(U\),
returns a learning signal to \(M\), and \(M\) uses that signal to
improve \((F, K)\) for the next action. Intelligence is not visible in
any single step; it is the rate at which the steps improve.

\subsubsection{Conclusion}\label{conclusion}

The formal results establish a minimal structural account of
intelligence under bounded context. The account rests on four
primitives: a world state space \(W\) generating uncertainty, a bounded
active context \(C_n\), a survival threshold \(U_{lethal}^d\), and a
growing reachable space \(\mathcal{T}_{reachable}\). From these:

\begin{itemize}
\tightlist
\item
  Knowledge and action are inseparable: \(K\) is constitutively indexed
  by \(F\), with an asymmetric retention criterion derived from survival
  pressure (Theorem 1).
\item
  Graph-structured \(K\) is necessary: bounded context requires
  progressive disclosure, which requires sub-linear retrieval, which
  requires stored adjacency --- a graph ({[}McCarthy2026a{]}).
\item
  Escalation dominates top-down allocation: bottom-up failure signaling
  is \(O(1)\) in steady state; top-down precision allocation degrades as
  \(O(|F|/C_n)\) (Theorem 2a).
\item
  The encoding hierarchy is a necessary consequence of bounded \(C_n\):
  proven knowledge promotes to lower levels, freeing the active context
  for genuine novelty (Theorem 4).
\item
  Agency \(\mathcal{A} > 0\) and learning efficiency \(\eta_M > 0\) are
  structural invariants of any persisting entity under selection
  pressure (Theorem 6).
\item
  Intelligence \(I(E) = \mathcal{A} \cdot \eta_M\) is the driven
  capacity to improve --- not what has been learned, but how efficiently
  evidence converts into better \((F, K)\).
\end{itemize}

The framework is substrate-independent: it specifies structural
requirements, not implementation. Any system satisfying these
constraints --- biological, artificial, or hybrid --- is a member of the
structural class. The biological architecture is accommodated as one
instance: heritable encoding at \(L_0\), individual learning at
\(L_n\), and cultural transmission as external \(K\) when the heritable
channel cannot carry individual \((K, F)\).

The theory makes testable predictions: that graph structure emerges
under bounded-context optimization (confirmed in the engineering system
that motivated the framework); that top-down orchestration degrades
under scaling (observable in both software and cognitive architectures);
and that the encoding hierarchy follows from cost structure rather than
requiring separate explanation.

Extensions to collective intelligence, superintelligence, and sentience
--- including formal definitions and convergence results --- are
developed in a companion paper.

\begin{center}\rule{0.5\linewidth}{0.5pt}\end{center}

\subsection{9. Related Work}\label{related-work}

\subsubsection{Graph Necessity and Autonomous
Learning}\label{graph-necessity-and-autonomous-learning}

The graph necessity result on which this paper builds is established in
{[}McCarthy2026a{]}: any system accumulating propositions under bounded
active context must maintain a directed graph with typed edges. The
proof proceeds from two independent arguments --- information
preservation (factoring is the only lossless space-creating operation,
and factoring is graph construction) and retention dynamics (sub-linear
dependency queries require stored adjacency). That paper also shows that
continuous-weight retrieval mechanisms (e.g., transformer self-attention)
implement graph structure rather than constituting an alternative to it.

Dupoux, LeCun, and Malik {[}DupouxLeCunMalik2026{]} propose a
dual-system architecture (observation + action) with meta-control
routing, motivated by cognitive science. Their framework is
architectural and prescriptive: it argues that autonomous learning
\emph{should} integrate observation, action, and meta-control. The
present work is complementary: rather than proposing what intelligent
systems should contain, we derive what they \emph{must} contain. Their
Systems A and B correspond to \(F\) (informed actions on \(W\)) and
\(K\) (observational propositions about \(W\)); their System M
corresponds to the meta-control aspect of \(M\) in this framework. The
structural results here --- type irreducibility of \(F\) and \(M\)
(Theorem 5), the necessity of \(\mathcal{A}\) and \(\eta_M\) (Theorem
6), and escalation dominance (Theorem 2a) --- provide formal grounding
for architectural choices that {[}DupouxLeCunMalik2026{]} motivates
empirically.

\subsubsection{Free Energy Principle and Active
Inference}\label{free-energy-principle-and-active-inference}

The Free Energy Principle {[}Friston2010{]} establishes that any
self-organizing system that maintains its existence necessarily
minimizes variational free energy --- formally equivalent to bounding
uncertainty about hidden world states. Active inference
{[}Friston2017{]} extends this to action selection: the organism selects
actions that minimize expected free energy, combining epistemic value
(information gain) and pragmatic value (goal proximity).

\textbf{Containment, not contradiction.} This framework does not
contradict FEP. FEP agents are members of this framework's structural
class: they satisfy \(\mathcal{A} > 0\) (self-organizing systems act),
\(\eta_M > 0\) (generative model improves over time), bounded \(C_n\)
(inference under resource limits), and a survival condition: maintaining
a Markov blanket (statistical separation of internal from external
states) is the formal characterization of what it means for an entity to
persist --- the same structural requirement as \(U < U_{lethal}^d\)
stated in the language of conditional independence rather than
uncertainty magnitude. Both describe maintained boundary; neither
implies the other's formalism, but they pick out the same class of
persisting systems. The variational free energy \(\mathcal{F}\)
decomposes into accuracy and complexity terms; the accuracy term
\(-\mathbb{E}_q[\log P(o \mid m, \pi)]\) is precisely the expected
surprisal \(H(W \mid K)\) of Definition 5. FEP agents therefore
instantiate this framework with \(U(w,K)\) realized as variational free
energy and \(M\) realized as Bayesian belief updating. The claim is not
that FEP is wrong. The claim is that it is a special case within a
broader structural class --- in the same way that Newtonian mechanics is
not wrong but is contained within the class of systems satisfying more
general symmetry constraints.

What the framework adds is four things FEP does not provide. First, FEP
derives the minimization behavior from the requirement of
self-organization but does not establish the minimization
\emph{imperative} --- why \(\mathcal{A} > 0\) is structurally necessary
for persistence. Theorem 6 establishes this as a necessary condition:
any entity that does not reduce uncertainty eventually fails and is
removed, so \(\mathcal{A} > 0\) is a structural invariant of persisting
systems. (The origin of \(\mathcal{A}\) --- how self-maintaining
structures bootstrap from initial conditions --- is a separate open
question; see Section 11.) Second, FEP's
generative model is unconstrained in structure; we derive a structural
constraint from survival pressure: \(K\) is necessarily indexed by
\(F\), with an asymmetric retention criterion that makes specific
testable predictions about what knowledge is retained and what is
pruned. Third, FEP's precision allocation model is top-down: higher
levels direct attention to lower levels by weighting prediction errors.
This is one specific implementation of \(M\). We show formally (Theorem
2a) that this implementation degrades inference quality as capability
grows under bounded context, because delegation overhead fills \(C_n\)
proportionally to \(|F|\). The escalation architecture --- bottom-up
failure signaling with promoted crystallization --- is a different
implementation of \(M\) that is \(O(1)\) in steady state. This is a
critique of FEP's architectural \emph{prescription} under bounded
context, not of its minimization objective or its descriptive
correctness. FEP's precision weighting may be optimal when \(C_n\) is
not the binding constraint; under bounded context, escalation dominates.

\subsubsection{Information-Theoretic
Accounts}\label{information-theoretic-accounts}

Shannon's mutual information {[}Shannon1948{]} grounds Definition 5
(\(U = H(w \mid K)\)) and Definition 3b
(\(G(f, K) = I(f\,;\,w \mid K)\)). The data processing inequality
{[}Cover2006{]} grounds the \(G \geq 0\) bound. The information
bottleneck principle {[}Tishby2011{]} treats the retention threshold for
knowledge-action coupling as a design parameter; this paper derives that
threshold from survival pressure, making it a consequence rather than a
choice. Schmidhuber's formal theory of intrinsic motivation
{[}Schmidhuber2010{]} formalizes curiosity as compression progress ---
the rate at which a learning algorithm improves its world model ---
which is structurally \(\eta_M\). The difference: Schmidhuber treats the
drive as a design objective; here \(\mathcal{A} > 0\) is established as
a necessary condition of persistence (Theorem 6), and \(\eta_M\) is
separated from it as a formally independent quantity.

\subsubsection{Cognitive Architectures}\label{cognitive-architectures}

Global Workspace Theory {[}Baars1988, Dehaene2011{]} holds that local
processors handle routine computation; the global workspace broadcasts
only when local processors fail. This is the escalation architecture in
cognitive science language. The contribution here: the escalation
criterion is derived from \(U\)-reduction cost (Theorem 2) rather than
assumed architecturally, the hierarchy is continuous rather than binary,
and Theorem 2a shows why top-down allocation from the global workspace
is inefficient. Kahneman's System 1/System 2 distinction
{[}Kahneman2011{]} is the popular form of the same binary --- corrected
here by the continuous \(\mathcal{E}(c, r)\) parameterization of
Definition 8. Gibson's affordances {[}Gibson1979{]} are action
possibilities indexed by organism capacities --- \(K\) indexed by \(F\)
under a different name. Theorem 1 provides what Gibson's account lacks:
a formal selection criterion for which affordances are retained.
Anderson and Schooler {[}AndersonSchooler1991{]} show empirically that
human memory retention mirrors the statistical structure of the
environment --- precisely the prediction of Theorem 1's retention
criterion applied to biological \(K\).

\subsubsection{Formal Theories of
Intelligence}\label{formal-theories-of-intelligence}

Hutter's AIXI {[}Hutter2005{]} defines the optimal agent as one that
maximizes expected reward over all computable environments, using
Solomonoff induction for prediction. Legg and Hutter
{[}LeggHutter2007{]} derive from this a formal definition of
intelligence as average performance across all computable environments
weighted by Kolmogorov complexity. The present framework differs in
three respects. First, AIXI assumes unbounded computation; the central
constraint here is bounded \(C_n\) and the consequences that follow from
it. Second, AIXI defines intelligence as performance; this framework
defines it as the rate of improvement (\(\eta_M\)), which is independent
of current performance level. Third, AIXI's environment class is all
computable functions; here the environment is physical (\(W\) with
positive entropy), and the survival threshold \(U_{lethal}^d\) is a
structural feature of the entity-world coupling, not a design parameter.

Chollet {[}Chollet2019{]} defines intelligence as skill-acquisition
efficiency over novel tasks --- structurally close to \(\eta_M\) as
defined in Definition 11a. The distinction: Chollet treats the prior
knowledge and experience of the agent as a confound to be controlled;
here prior \(K\) is constitutive of intelligence because \(\eta_M\)
operates on \((F, K)\). Chollet's ``scope'' and ``generalization
difficulty'' map to the growth of \(\mathcal{T}_{reachable}\); his
``priors'' correspond to the encoding hierarchy's lower levels. The
frameworks are compatible but ask different questions: Chollet asks how
to measure; this paper asks what structural conditions are necessary.

\subsubsection{Bounded Rationality and Resource-Rational
Analysis}\label{bounded-rationality}

Simon {[}Simon1955{]} established that rational agents operating under
computational and informational constraints do not optimize globally but
satisfice --- selecting the first option that exceeds a threshold. The
bounded context constraint \(C_n\) formalizes one specific source of the
bounds Simon identified. The retention criterion (Theorem 1) is a
satisficing criterion: retain \(p\) if benefit is not certainly zero,
rather than computing the globally optimal \(K\).

Griffiths, Lieder, and Goodman {[}GriffithsEtAl2015{]} formalize
resource rationality: the optimal algorithm given computational
constraints may differ from the optimal algorithm given unlimited
resources. This is precisely the argument of Theorem 2a --- top-down
precision allocation is optimal without resource constraints; under
bounded \(C_n\), escalation dominates. The encoding hierarchy
(Definition 8) is a resource-rational architecture: it allocates costly
active inference only to genuinely uncertain problems, handling resolved
problems at lower cost.

Clark {[}Clark2015{]} synthesizes predictive processing with embodied
cognition, arguing that the brain is fundamentally a prediction engine
that minimizes prediction error hierarchically. This is the cognitive
science framing of what Theorem 2 formalizes: hierarchical error
minimization under bounded context. Clark's ``action-oriented predictive
processing'' corresponds to the \(K\)-indexed-by-\(F\) structure of
Theorem 1 --- predictions are constitutively tied to actions.

Pearl {[}Pearl2009{]} provides the formal framework for causal reasoning
in directed graphs. The knowledge graph \(K\) as defined here (Definition
2) encodes inferential and causal relationships as a directed graph; the
Bayesian network interpretation of Definition 5 is grounded in Pearl's
framework. Pearl's do-calculus formalizes interventional reasoning ---
the distinction between observing and acting --- which corresponds to
the distinction between \(K\) (what is known) and \(F\) (what can be
done). The retention criterion of Theorem 1 can be read as a
causal relevance filter: retain \(p\) if there exists a causal path
from \(p\) to some \(f \in F\) that reduces \(U\).

\begin{center}\rule{0.5\linewidth}{0.5pt}\end{center}

\subsection{10. Discussion}\label{discussion}

The formal results have implications beyond the immediate theorems.

\textbf{The encoding principle.} Corollary 6
implies that as \(K\) grows, informed actions are promoted
down the encoding hierarchy. Work that is currently active reasoning
(uncertain, costly, engaged at \(L_n\)) becomes reflex, then encoding,
then crystallized. By Theorem 2 and
Definition 8, this frees the active context window for the next genuine
frontier of uncertainty. \(\mathcal{A}\) does not converge when its
current domain is solved --- it reaches outward.

\textbf{The civilization-scale knowledge system as empirical evidence.}
An internal-model account predicts decelerating knowledge growth: larger
models contain more variance, \(\rho\) degrades, and recovery requires
expanding capacity (Theorem 2a). Civilizational knowledge growth has
accelerated. The explanation: intelligence externalized \(K\) into
graph-structured repositories and retrieves \(K^{[f]}\) into active
context per problem. This is the escalation architecture (Theorem 2) and
graph necessity result ({[}McCarthy2026a{]}) at civilizational scale. Two
predictions an internal-model account cannot make: (i) knowledge growth
accelerates as \(K\) grows (Corollary 6 and Theorem 4: proven actions
promote to lower levels, freeing \(L_n\) for new learning), and (ii) the
external
architecture self-assembles as a graph with provenance. Both are
observed.

\textbf{Limits of the formal account.} The theory establishes structural
requirements for individual intelligence under bounded context. It does
not address collective intelligence, convergence to general intelligence,
or the conditions under which shared knowledge reduces uncertainty faster
than any individual. These extensions are developed in a companion paper.

\begin{center}\rule{0.5\linewidth}{0.5pt}\end{center}

\subsection{11. Open Questions}\label{open-questions}

\begin{enumerate}
\def\labelenumi{\arabic{enumi}.}
\item
  \textbf{The origin of \(\mathcal{A}\) (distinct from its necessity).}
  Theorem 6 establishes that \(\mathcal{A} > 0\) is a necessary
  condition for persistence: it is a structural invariant of any entity
  that survives under selection pressure with bounded context. This is a
  selection argument --- it explains why \(\mathcal{A} > 0\) entities
  predominate, not how \(\mathcal{A} > 0\) structures arise from initial
  conditions where no such structures exist. The origin question is
  logically independent of the necessity result and remains open.

  The question: what mechanism produces the first self-maintaining
  structures with primitive \(\mathcal{A} > 0\) from \(\mathcal{A} = 0\)
  initial conditions? Can \(\mathcal{A}\) be constructed, or only
  selected for once it exists?

  \emph{A conjectured answer (not a formal result).} \(\mathcal{A}\) may
  be thermodynamically instantiated. \(U(w, K) = H(w \mid K)\) is
  information entropy. Reducing \(U\) is reducing entropy locally. In
  any universe with entropy gradients, structures that maintain
  themselves against thermodynamic pressure persist preferentially. The
  first self-maintaining structure was, in the minimal sense, bounding
  its own uncertainty: \(\mathcal{A}\) in primitive form. Not designed.
  A consequence of the physics of a world that destroys things.

  The deeper question is therefore: is agency a natural consequence of
  the universe? The conjectured answer is yes: \(\mathcal{A}\) is the
  thermodynamic attractor of any entropy-gradient-rich world given
  sufficient time. Wherever entropy gradients exist and structures can
  form, the structures that persist are those that resist dissolution by
  reducing local uncertainty. Given enough time, those structures
  compound. Intelligence is what that compounding produces.

  Formal derivation of this claim from thermodynamic first principles
  (connecting \(U_{lethal}^d\) to physical entropy thresholds and
  demonstrating the inevitability of primitive \(\mathcal{A}\) from
  dissipative system dynamics) remains open. See {[}Prigogine1984{]}
  (dissipative structures), {[}England2013{]} (dissipative adaptation),
  and {[}Schrodinger1944{]} (negative entropy and life) for adjacent
  physical frameworks.
\item
  \textbf{The direction of \(\mathcal{A}\).} Is agency a uniform drive
  (a general orientation toward uncertainty reduction), or does it carry
  direction? Survival pressure selects for reducing \emph{specific}
  uncertainties. Does \(\mathcal{A}\) encode this specificity, or does
  \(F\) and \(K\) direct it toward particular regions of \(W\) after the
  fact?
\item
  \textbf{Perception and \(K\) construction.} How does \(\sigma\)
  determine which features of \(W_{obs}\) are encoded into \(K\)? Is
  \(\sigma\) itself subject to gradient descent, and if so, at what
  level?
\item
  \textbf{Bound composition.} Can bounds compose, i.e., can
  \(f_1 \circ f_2\) produce informed actions covering regions of \(W\)
  neither \(f_1\) nor \(f_2\) could cover alone? What are the rules of
  composition?
\item
  \textbf{The novelty boundary (theory limit).} The framework defines
  informed action over regions of \(W\) where \(K\) provides grounding.
  At \(U(w, K) \approx 1\) --- where no existing informed action applies
  --- the theory does not specify \(E\)'s behavior. This is a boundary
  condition, not a design choice: the model is silent on uninformed
  action by construction. Behavior in this regime (random exploration,
  interpolation, escalation to \(M\)) is outside the formal scope and
  would require extension of the framework.
\item
  \textbf{Competing bounds.} When multiple bounds \(f_i \in F\) are
  applicable to the same region of \(W\), how does \(E\) select among
  them? Is there a selection function, and is it itself in \(F\) or
  \(M\)?
\item
  \textbf{The \(M^{(n)}\) hierarchy.} The framework defines \(M\) as
  operating on \((F, K)\) (Definition 11) and notes that the same type
  applies recursively. The formal questions are: how many levels are
  required, what determines when another level becomes necessary, and
  whether the hierarchy has a natural termination point?

  The framework's answer to the first question is empirical, not
  prescribed: one level of \(M\) suffices until it does not. As long as
  the gradient on \((F, K)\) can be followed by improving \((F, K)\)
  directly, \(M^{(1)}\) is adequate. When it cannot --- when the
  learning mechanism itself is the bottleneck --- selection pressure
  propagates to \(M^{(2)}\). A single level is fine, until it isn't. The
  required depth is determined by the structure of the problem and the
  selection pressure it generates, not by the theory in advance. The
  author's computational system has two explicit levels; there is no
  formal argument that two is sufficient in general.

  \emph{Hypothesis (M-depth driven by selection pressure).} Each level
  \(M^{(n)}\) is subject to the same minimization imperative as the
  level below: it will be displaced or refined if a better \(M^{(n+1)}\)
  exists and the selection pressure is sufficient to instantiate it. The
  hierarchy has a natural termination at level \(n^*\) only if the
  gradient at that level reaches zero --- no available \(M^{(n^*+1)}\)
  can improve the learning mechanism further given current \(W\).
  Whether such a fixed point exists, and whether it is reachable in
  finite time under realistic selection pressure, is open. Evolution
  {[}Hinton1987{]} and meta-learning {[}FinnEtAl2017, Thrun1998{]} are
  empirical candidates for \(M^{(2)}\); whether anything instantiates
  \(M^{(3)}\) in known systems remains unclear.
\item
  \textbf{Genetic encoding and the developmental sufficiency
  requirement.} The encoding hierarchy (Definition 8) describes a
  ratchet: proven behaviors promote toward \(L_0\) as they meet the
  rigor threshold \(\theta_0 \approx 1\) (Theorem 4). The question is
  what the asymptotic product of evolutionary gradient descent encodes
  {[}Schrodinger1944, MaynardSmith1995, Hinton1987, England2013{]}.

  The framework does not resolve the content of genetic encoding ---
  what specific mixture of \(M\), \((F_0, K_0)\), and developmental
  constraint DNA carries across all levels is an open empirical
  question. The claim is narrower: whatever genetic encoding carries
  must be \emph{sufficient to develop the entity} --- sufficient to
  bootstrap an organism capable of building \((F, K)\) through
  interaction with \(W\). The model is compatible with genetic encoding
  at any level of the hierarchy, in any proportion, provided the
  developmental output is a system that can run the gradient.

  What the framework does predict about the evolutionary rigor
  threshold: specific \((F, K)\) are environment-dependent --- a skill
  that reduces \(U\) in one environment may not transfer across the full
  distribution of environments the lineage will encounter. Content that
  clears \(\theta_0\) must generalize across that distribution.
  Developmental machinery --- the programs, architectures, and gradient
  signals that build \((F, K)\) --- is more likely to generalize than
  specific content, because it is environment-agnostic. The framework
  predicts selection pressure favoring developmental generality over
  content specificity at the highest rigor levels; it does not predict
  that specific \((F_0, K_0)\) is absent.

  This is consistent with the observation that we are not born with the
  context window of our parents: individual \((K, F)\) --- everything
  learned in one lifetime --- clears \(\theta_n\) but not \(\theta_0\).
  The inheritance mechanism transmits what is sufficient for
  development, not what was learned. The diversity of knowledge across
  individuals with similar genetic endowment is the expected outcome:
  the same developmental machinery operating on different environments
  produces different \((K, F)\).

  Cultural transmission (language, writing, institutions) provides a
  parallel channel for \((K, F)\) at lower fidelity but higher speed ---
  shared \(K\) functioning as external inheritance,
  serving a transmission function that heritable encoding cannot match for
  individual-lifetime content. The two channels operate at different
  rigor thresholds and different timescales.

  Formal questions that remain open: the precise allocation between
  \(M\)-encoding and \((F_0, K_0)\)-encoding across evolutionary
  lineages; whether the modularity of genetic regulatory networks is
  derivable from the type structure of
  \(M : (F, K) \times \mathcal{L} \rightarrow (F, K)\) at the
  evolutionary scale; and what determines \(\theta_0\) in terms of
  environmental distribution statistics.
\end{enumerate}

Further open questions concerning collective intelligence, the AGI
condition, and the structure of genetic encoding are developed in a
companion paper.

\end{document}
